\documentclass[11pt,a4paper]{article}

\usepackage[T1]{fontenc}
\usepackage[utf8]{inputenc}
\usepackage[british]{babel}
\usepackage{lmodern}
\usepackage{microtype}
\usepackage[a4paper,margin=1in]{geometry}
\usepackage{amsmath,amssymb,amsthm,mathtools}
\usepackage{enumitem}
\usepackage{xcolor}
\usepackage{hyperref}
\hypersetup{
  colorlinks=true,
  linkcolor=blue!60!black,
  urlcolor=blue!60!black,
  hypertexnames=false,
  pdftitle={Trace-Tempered Choice: Lean-Faithful Appendices to Theorem 1},
  pdfauthor={Daniele Condorelli}
}
\setlength{\emergencystretch}{2em}
\interfootnotelinepenalty=10000

\newtheorem{theorem}{Theorem}
\newtheorem{proposition}[theorem]{Proposition}
\newtheorem{lemma}[theorem]{Lemma}
\newtheorem{corollary}[theorem]{Corollary}
\theoremstyle{definition}
\newtheorem{definition}[theorem]{Definition}
\newtheorem{remark}[theorem]{Remark}

\newcommand{\D}{\Delta}
\newcommand{\E}{\mathbb E}
\newcommand{\I}{I}
\newcommand{\Sh}{H}
\newcommand{\supp}{\operatorname{supp}}
\newcommand{\Id}{\operatorname{Id}}
\newcommand{\R}{\mathbb R}
\newcommand{\Rnum}{\mathbb R}
\newcommand{\KL}{D_{\mathrm{KL}}}
\newcommand{\1}{\mathbf 1}

% The main paper defines and labels these two transformations before the
% appendix.  The preview supplies compact copies so that the fragment compiles
% and reads on its own.
\newcommand{\recordprocessingtag}{A.1}
\newcommand{\actionprocessingtag}{A.2}
\newcommand{\eqrefrecordprocessing}{\eqref{eq:record-processing}}
\newcommand{\eqrefactionprocessing}{\eqref{eq:action-processing}}

\begin{document}

\begin{center}
{\Large\bfseries Trace-Tempered Choice}\par
\vspace{0.4em}
{\large Lean-faithful appendices to Theorem 1}\par
\vspace{0.6em}
Daniele Condorelli
\end{center}

\paragraph{Standing notation.}
All payoff, action, and record alphabets are finite, and action and record
alphabets are nonempty.  The primitive relations, common-payoff block
environments, and Axioms \textup{(A1)--(A8)} are those in the paper.  A
payoff-preserving record processor $T$ and an action report $S$ satisfy,
respectively,
\begin{align}
(KT)(o,r'\mid a)
  &=\sum_r K(o,r\mid a)T(r'\mid o,r),
  \label{eq:record-processing}\\
(qS)(b)\widehat K(o,r\mid b)
  &=\sum_a q(a)S(b\mid a)K(o,r\mid a).
  \label{eq:action-processing}
\end{align}
The second identity is imposed without division and hence leaves rows at null
reported actions unrestricted.

\tableofcontents

\appendix
\input{theorem1_lean_faithful_appendix}

\clearpage
% Generated from the leanpurestatement and leanpureproof regions of
% trace_tempered_choice_v3.tex.  Run extract_lean_faithful_pure_trace.sh
% after editing either marked region; do not edit this file by hand.

\section{The pure-trace characterization}
\label{pt:sec:environment}\label{lf:app:pure-trace}
% ===============================================================

Let $A$ be a finite set of \emph{actions} and $X$ a finite set of \emph{records}.  A channel is a stochastic kernel
$P:A\to\Delta(X)$.  For each finite channel $P$, the primitive behavioural datum is a binary relation
\[
\succeq_P\quad\text{on}\quad \Delta(A).
\]
Condition \textup{(P1)} below requires this relation to be a weak order.  The interpretation is that
$q\succeq_P q'$ means: in the fixed environment $P$, lottery $q$ displays at least as much
traceability as $q'$.
The environment $P$ is not an object of choice.  A lottery is the procedure being evaluated; the channel is the fixed
observation technology through which the realised act may be recorded by records.

\paragraph{Comparison environments.}
To compare lotteries carried by different channels without abandoning channel-indexed preferences, place the channels in
labelled, disjoint blocks of one larger environment.  For channels $P:A\to\Delta(X)$ and $Q:B\to\Delta(Y)$, write
$P\sqcup Q$ for the block-diagonal channel with action alphabet
$(A\times\{0\})\cup(B\times\{1\})$ and record alphabet
$(X\times\{0\})\cup(Y\times\{1\})$.  It sends $(a,0)$ to $(x,0)$ with probability $P(x\mid a)$ and $(b,1)$ to $(y,1)$
with probability $Q(y\mid b)$; all cross-block probabilities are zero.  For $q\in\Delta(A)$ and
$p\in\Delta(B)$, define lotteries on the tagged action alphabet by
\[
\begin{aligned}
q^0(a,0)&:=q(a), & q^0(b,1)&:=0,\\
p^1(a,0)&:=0,    & p^1(b,1)&:=p(b),
\end{aligned}
\qquad (a\in A,\ b\in B).
\]
The superscript is a block label, not a power.  Accordingly, $q^0\succeq_{P\sqcup Q}p^1$ is an ordinary preference
comparison inside one fixed environment.\footnote{For a complete numerical example, let
$A=\{a_1,a_2\}$, $B=\{b_1,b_2\}$, $X=\{x_0,x_1\}$, and $Y=\{y_0,y_1\}$.  With rows indexed by actions and columns
by records, take
$P=\bigl(\begin{smallmatrix}0.1&0.9\\0.8&0.2\end{smallmatrix}\bigr)$ and
$Q=\bigl(\begin{smallmatrix}0.4&0.6\\0.6&0.4\end{smallmatrix}\bigr)$.
Order the tagged actions as $(a_1,0),(a_2,0),(b_1,1),(b_2,1)$ and the tagged records as
$(x_0,0),(x_1,0),(y_0,1),(y_1,1)$.  Then $P\sqcup Q=\operatorname{diag}(P,Q)$: the first two actions generate only
the first two records, and the last two actions only the last two.  If $q=(0.3,0.7)$ and $p=(0.8,0.2)$, their
embeddings are $q^0=(0.3,0.7\,;\,0,0)$ and $p^1=(0,0\,;\,0.8,0.2)$, where the semicolon separates the two blocks.  Thus
$q^0\succeq_{P\sqcup Q}p^1$ says precisely that lottery $q$, observed through $P$, is at least as traceable as lottery
$p$, observed through $Q$; it does not treat either channel as an object of choice.  Under the representation proved
below, the two sides yield approximately $0.330$ and $0.019$ bits of mutual information, respectively, and hence the
comparison is strict in this example.}
When the two blocks have the same action alphabet, $q^0$ and $q^1$ are the two tagged copies of the same lottery $q$. For a finite family
$\{P_i:A_i\to\Delta(X_i)\}_{i\in K}$, write $\bigsqcup_{i\in K}P_i$ for the analogous finite block environment and
$q_i^i$ for the lottery that assigns $q_i(a)$ to each tagged action $(a,i)$ and zero to every other block.

\paragraph{Pure-record conditions.}
The conditions are imposed directly on the family $\{\succeq_P\}_P$.
Five conditions are imposed on the family. The first three are structural:
      each environment carries a continuous weak order with a fixed orientation
      (P1, P2), and labelled blocks provide the comparison framework (P3).
      The substantive restrictions are two: garbling cannot increase
      traceability (P4), and sequential observation is evaluated branch by
      branch (P5).


\begin{description}[leftmargin=*,style=sameline]

\item[(P1) Weak order and local non-triviality.]
For every finite channel $P:A\to\Delta(X)$, $\succeq_P$ is complete and transitive on $\Delta(A)$.  The family is
locally non-trivial: for every finite action set $A$ with $|A|\ge 2$ and every full-support $q\in\Delta(A)$,
\[
q^0\succ_{\Id_A\sqcup U_A}q^1,
\]
where $\Id_A:A\to\Delta(A)$ is full revelation and $U_A:A\to\Delta(\{\ast\})$ is the one-record uninformative channel.

This condition gives each observation environment an ordinary weak order over act lotteries and fixes the orientation of the
comparison.  When the realised act is genuinely uncertain, observing it directly is strictly more traceable than observing
nothing about it.

\item[(P2) Continuity.]
For every pair of finite alphabets $(A,X)$, the set
\[
\{(P,q,q')\in \Delta(X)^A\times\Delta(A)\times\Delta(A):q\succeq_P q'\}
\]
is closed in the usual Euclidean product topology.  Thus within-environment rankings vary continuously with both the
lotteries and the fixed channel.

Continuity says that traceability comparisons are stable under small changes in the procedure and the environment.  Small
modelling perturbations should not create abrupt reversals.  The condition is stated entirely in the primitives: it quantifies
over channels and lotteries, not over derived Bayesian objects.  Continuity of block comparisons in posterior-law
convergence at a fixed full-support prior is not assumed; it is derived below
(Lemma~\ref{pt:lem:plcont}).

\item[(P3) Block-comparison coherence.]
\emph{Duplication:} for every channel $P:A\to\Delta(X)$ and every $q,q'\in\Delta(A)$,
\[
q\succeq_P q'
\quad\Longleftrightarrow\quad
q^0\succeq_{P\sqcup P}(q')^1.
\]
\emph{Irrelevant blocks:} for every finite block environment
$\bigsqcup_{k\in K}P_k$, every ordered pair of distinct blocks $i,j\in K$, and every
$q_i\in\Delta(A_i)$ and $q_j\in\Delta(A_j)$,
\[
q_i^i\succeq_{\bigsqcup_{k\in K}P_k}q_j^j
\quad\Longleftrightarrow\quad
q_i^0\succeq_{P_i\sqcup P_j}q_j^1.
\]
Applied at the ordered pair $(1,0)$ of a two-block environment, the second clause exchanges block labels; this is
how reversed block comparisons are brought to canonical form throughout
(Lemma~\ref{pt:lem:blockcoh}\textup{(a)}).

Labelled blocks are a way to place procedures carried by different channels in a common comparison problem.  The labels
should not themselves create traceability, alter the ranking inside a duplicated environment, or let unused blocks affect
the comparison between active blocks.

\item[(P4) Split data-processing aversion.]
Condition \textup{(P4)} is the conjunction of the following two clauses.  This
is the formulation used by the formal theorem: the record and action
coordinates are processed separately, and no convention is imposed on null
reported actions.

\emph{\textup{(DP-rec)} Record post-processing.}  For every channel
$P:A\to\Delta(X)$, every stochastic $T:X\to\Delta(X')$, and every
$q\in\Delta(A)$, define
\[
(PT)(x'\mid a):=\sum_{x\in X}P(x\mid a)T(x'\mid x).
\]
Then
\[
q^0\succeq_{P\sqcup PT}q^1.
\]

\emph{\textup{(DP-act)} Bayesian action coarsening.}  For every
$P:A\to\Delta(X)$, every stochastic $S:A\to\Delta(A')$, every
$q\in\Delta(A)$, write
\[
(qS)(a'):=\sum_{a\in A}q(a)S(a'\mid a).
\]
For every channel $\widehat P:A'\to\Delta(X)$ satisfying
\begin{equation}
(qS)(a')\widehat P(x\mid a')
=\sum_{a\in A}q(a)P(x\mid a)S(a'\mid a)
\quad\text{for every }(a',x)\in A'\times X,
\tag{DP-act-law}\label{pt:eq:dpact}
\end{equation}
we have
\[
q^0\succeq_{P\sqcup \widehat P}(qS)^1.
\]
When $qS$ has full support, \eqref{pt:eq:dpact} determines $\widehat P$
uniquely; otherwise the clause applies to every channel consistent with it.
In particular, bijective relabellings are neutral by applying
\textup{(DP-act)} in both directions.

The act whose trace is evaluated may be reported through a possibly noisy description, and the visible record may be
garbled after it is generated.  Neither operation can create observable trace: the original procedure is weakly more
traceable than any procedure obtained by either operation.

\begin{remark}\label{pt:rem:dpform}
The two primitive clauses imply the convenient simultaneous form used in a
few calculations.  Let $S:A\to\Delta(A')$ and $T:X\to\Delta(X')$, and let
$P':A'\to\Delta(X')$ satisfy
\begin{equation}
(qS)(a')\,P'(x'\mid a')
=
\sum_{a\in A}\sum_{x\in X}q(a)P(x\mid a)S(a'\mid a)T(x'\mid x)
\quad\text{for every }(a',x')\in A'\times X'.
\tag{DP}\label{pt:eq:dpjoint}
\end{equation}
First apply \textup{(DP-rec)} to form $PT$, and then apply
\textup{(DP-act)} to $(q,PT)$ through $S$.  Equation~\eqref{pt:eq:dpjoint}
says exactly that the supplied $P'$ is an allowed Bayesian completion of this
second step, including on null reported actions.  Chain the two comparisons in
a common three-block environment using Conditions \textup{(P1)} and
\textup{(P3)}.  This yields
$q^0\succeq_{P\sqcup P'}(qS)^1$.  Conversely, the simultaneous form
specialises to each named clause.  Thus the two presentations are equivalent
under \textup{(P1)} and \textup{(P3)}, but the proposition below takes the two
named clauses as its literal premise.  If $(qS)(a')=0$, both
\eqref{pt:eq:dpact} and \eqref{pt:eq:dpjoint} leave the corresponding channel
row unrestricted; null reports require no convention.
\end{remark}

\end{description}

\paragraph{Reached records and posteriors.}
For a lottery $q\in\Delta(A)$ and a channel $P:A\to\Delta(X)$, write
\[
m_{q,P}(x):=\sum_{a\in A}q(a)P(x\mid a),
\qquad
r_x^{\,q,P}(a):=\frac{q(a)P(x\mid a)}{m_{q,P}(x)}
\quad\text{if }m_{q,P}(x)>0
\]
for, respectively, the probability of record $x$ and the posterior belief about the realised action.  Call $x$ reached
when $m_{q,P}(x)>0$; all branch comparisons below are restricted to reached records.

\begin{description}[leftmargin=*,style=sameline]

\item[(P5) Branchwise continuation monotonicity.]
Let $P_1:A\to\Delta(X_1)$ be a first-stage channel and, after each first-stage record $x$, let
$Q^x,R^x:A\to\Delta(X_2^x)$ be two alternative continuation channels.  Write
$P_1\triangleright\{Q^x\}$, and analogously $P_1\triangleright\{R^x\}$, for the compound channel with record alphabet
$\bigcup_{x\in X_1}(\{x\}\times X_2^x)$ that reports both stages:
\[
(P_1\triangleright\{Q^x\})(x,z\mid a):=P_1(x\mid a)Q^x(z\mid a).
\]
Fix $q\in\Delta(A)$ and let
$m(x):=m_{q,P_1}(x)$ and $r_x:=r_x^{\,q,P_1}$ for all branches with $m(x)>0$.

If
\[
r_x^0\succeq_{Q^x\sqcup R^x}r_x^1
\qquad\text{for every }x\text{ with }m(x)>0,
\]
then
\[
q^0
\succeq_{(P_1\triangleright\{Q^x\})\sqcup(P_1\triangleright\{R^x\})}
q^1.
\]
If at least one reached branch comparison is strict, then the aggregate comparison is strict.

Sequential observation is evaluated branch by branch.  After a first-stage signal, if one continuation leaves at least as
much trace as another at every posterior that can be reached, then replacing the continuation branch by branch should not
reduce overall traceability.  If the improvement is strict on a reached branch, the aggregate comparison is strict.

\end{description}

\begin{remark}\label{pt:rem:A5role}
Condition~\textup{(P5)} plays in this system the role the sure-thing principle plays in Savage's: it is a dominance
condition across the branches of a resolving observation, with zero-probability branches unconstrained exactly as
Savage's null events are.  It is also where all cardinal content resides.  The other conditions are ordinal,
environment-by-environment conditions; it is \textup{(P5)}, through its derived consequences---public-coin
independence (Lemma~\ref{pt:lem:publiccoin}), background inertness (Lemma~\ref{pt:lem:background}), and the chain rule
(Lemma~\ref{pt:lem:chain})---that aligns the initially separate scales across priors and forces branch-additive
aggregation.  The strictness clause is not decorative: it is what allows the converse directions of
Lemmas~\ref{pt:lem:publiccoin} and~\ref{pt:lem:background} to be extracted from a dominance condition.
For example, a mutual-information index capped at a fixed level satisfies
\textup{(P1)}--\textup{(P4)} but fails exactly \textup{(P5)}.  Compound observation environments are
accordingly the natural locus of experimental testing: they are where branch-additive and non-additive criteria come
apart.
\end{remark}

% ===============================================================
\subsection{Auxiliary result}
% ===============================================================

Throughout, logarithms are base two, so information is measured in bits; changing the base merely rescales all
information values.
For a finite distribution $x$, write
\[
\Sh(x):=-\sum_z x(z)\log x(z),
\qquad 0\log 0:=0,
\]
and, for a prior--channel pair $(q,P)$, define
\[
I_{q,P}(A;X)
:=
\Sh(m_{q,P})-\sum_{a\in A}q(a)\Sh\bigl(P(\cdot\mid a)\bigr).
\]
When the alphabets are clear, write $I_{q,P}$; write $I$ for the map $(q,P)\mapsto I_{q,P}$.

\begin{proposition}[Pure-trace characterisation]\label{pt:prop:main}
For a family $\{\succeq_P\}_P$ as described above, the following are equivalent.
\begin{enumerate}[label=\textup{(\roman*)},leftmargin=*]
\item The family satisfies Conditions~\textup{(P1)--(P5)}.
\item For every finite channel $P:A\to\Delta(X)$ and every $q,q'\in\Delta(A)$,
\[
q\succeq_P q'
\quad\Longleftrightarrow\quad
I_{q,P}(A;X)\ge I_{q',P}(A;X).
\]
\end{enumerate}
Moreover, under these equivalent conditions, block-supported cross-channel comparisons are represented on the same
scale: for every finite block environment $\bigsqcup_{k\in K}P_k$, every two distinct blocks $i,j\in K$, and every
$q_i\in\Delta(A_i)$, $q_j\in\Delta(A_j)$,
\[
q_i^i\succeq_{\bigsqcup_{k\in K}P_k}q_j^j
\quad\Longleftrightarrow\quad
I_{q_i,P_i}(A_i;X_i)\ge I_{q_j,P_j}(A_j;X_j).
\]
\end{proposition}

\subsection{Proof of the pure-trace characterization}\label{pt:app:sufficiency}
% ===============================================================

\subsubsection{Necessity}

For every finite channel $P:A\to\Delta(X)$, define
\[
q\succeq^I_P q'
\quad\Longleftrightarrow\quad
I_{q,P}(A;X)\ge I_{q',P}(A;X).
\]
We verify that the family $\{\succeq^I_P\}_P$ satisfies Conditions~\textup{(P1)--(P5)}: this is the implication
\textup{(ii)$\Rightarrow$(i)} of Proposition~\ref{pt:prop:main}, and it establishes the consistency of the condition
system.

The \emph{posterior law} induced by $(q,P)$ is the finite distribution
\[
\mu_{q,P}:=
\sum_{x:m_{q,P}(x)>0}m_{q,P}(x)\,\delta_{r_x^{\,q,P}}
\]
over posterior beliefs.  Bayes plausibility gives $\int r\,d\mu_{q,P}(r)=q$.  We repeatedly use the elementary identities
\[
I_{q,P}
=
\Sh(m_{q,P})-\sum_{a\in A}q(a)\Sh(P(\cdot\mid a)),
\tag{N1}
\]
and
\[
I_{q,P}
=
\Sh(q)-\int_{\Delta(A)}\Sh(r)\,d\mu_{q,P}(r).
\tag{N2}
\]
Both formulae are read on the support of the relevant distributions; zero-probability rows and records do not affect the
value.

\emph{Condition \textup{(P1)}.}
For each fixed $P$, the relation $\succeq^I_P$ is complete and transitive because it is represented by the real-valued
function $q\mapsto I_{q,P}$.  If $|A|\ge2$ and $q$ has full support, then full revelation gives
\[
I_{q,\Id_A}=\Sh(q)>0,
\]
whereas the one-record uninformative channel $U_A$ gives $I_{q,U_A}=0$.  Hence
\[
q^0\succ^I_{\Id_A\sqcup U_A}q^1.
\]

\emph{Condition \textup{(P2)}.}
For fixed finite alphabets, (N1) shows that $(P,q)\mapsto I_{q,P}$ is continuous, since Shannon entropy is continuous on
every finite simplex.  Therefore the set
\[
\{(P,q,q'):q\succeq^I_P q'\}
=
\{(P,q,q'):I_{q,P}\ge I_{q',P}\}
\]
is closed.  Posterior-law continuity of block comparisons is not an condition; for the mutual-information family it also
holds directly, by (N2) and continuity of $\Sh$ on the compact simplex $\Delta(A)$, consistently with its derivation
from the conditions in Lemma~\ref{pt:lem:plcont} below.

\emph{Condition \textup{(P3)}.}
For every channel $P$ and lotteries $q,q'\in\Delta(A)$,
\[
I_{q^0,P\sqcup P}=I_{q,P},
\qquad
I_{(q')^1,P\sqcup P}=I_{q',P},
\]
because a block-supported lottery puts probability one on its own block and the block label is constant under that
lottery.  Hence
\[
q\succeq^I_P q'
\quad\Longleftrightarrow\quad
q^0\succeq^I_{P\sqcup P}(q')^1.
\]
The same observation proves invariance to adding or deleting irrelevant blocks.  If $q_i^i$ and $q_j^j$ are supported on
blocks $i$ and $j$, their mutual-information values in $\bigsqcup_{k\in K}P_k$ are exactly
\[
I_{q_i,P_i}
\quad\text{and}\quad
I_{q_j,P_j},
\]
independently of all other blocks.

\emph{Condition \textup{(P4)}.}
It is useful to verify the stronger simultaneous inequality from
Remark~\ref{pt:rem:dpform}, which immediately supplies both named clauses.
Let $S:A\to\Delta(A')$ and $T:X\to\Delta(X')$ be stochastic, and let $P'$ satisfy \eqref{pt:eq:dpjoint}.  Generate
$A\sim q$; then, conditionally independently given $A$, generate $A'$ from $S(\cdot\mid A)$ and $X$ from
$P(\cdot\mid A)$; finally generate $X'$ from $T(\cdot\mid X)$.  The joint law of $(A,A',X,X')$ is
\[
q(a)S(a'\mid a)P(x\mid a)T(x'\mid x),
\]
so $A'$ is conditionally independent of $(X,X')$ given $A$, and $X'$ is conditionally independent of $(A,A')$ given $X$.
Hence
\[
A'\longleftarrow A\longrightarrow X\longrightarrow X'
\]
is a Markov chain in the sense that $A'\to A\to X\to X'$ satisfies the required conditional independences.  Data
processing applied to the fork $A'\leftarrow A\to X$ gives $I(A;X)\ge I(A';X)$, and data processing applied to
$A'\to X\to X'$ gives $I(A';X)\ge I(A';X')$.  Therefore
\[
I_{q,P}(A;X)\ge I(A';X').
\]
By construction the joint law of $(A',X')$ is exactly $(qS)(a')P'(x'\mid a')$, the left-hand side of
\eqref{pt:eq:dpjoint}, so $I(A';X')=I_{qS,P'}(A';X')$.  Rows $a'$ with $(qS)(a')=0$ carry no mass under $qS$, so this value
is the same for every $P'$ satisfying \eqref{pt:eq:dpjoint}, however the null rows are completed.  Using the block
computation of \textup{(P3)} above to evaluate the two blocks separately,
\[
q^0\succeq^I_{P\sqcup P'}(qS)^1.
\]
The first named clause is the case $S=\Id_A$ and gives
$I_{q,P}(A;X)\ge I_{q,PT}(A;X')$.  For later use, in the action-only case $T=\Id_X$ write $S^qP$ for the conditional
channel from the reported action to the record on every reached row:
\[
S^qP(x\mid a')
:=
\frac{\sum_{a\in A}q(a)S(a'\mid a)P(x\mid a)}{(qS)(a')}
\qquad\text{when }(qS)(a')>0.
\]
A \emph{completion} of $S^qP$ is any channel $\widehat P:A'\to\Delta(X)$ that agrees with this formula on reached rows;
its values on null rows are unrestricted.  The second named clause gives
$I_{q,P}(A;X)\ge I_{qS,\widehat P}(A';X)$ for every completion $\widehat P$ of $S^qP$.

\emph{Condition \textup{(P5)}.}
Fix a first-stage channel $P_1:A\to\Delta(X_1)$.  Let
\[
m(x)=m_{q,P_1}(x),
\qquad
r_x=r_x^{\,q,P_1}.
\]
For a continuation profile $\mathcal Q=\{Q^x\}_x$, the chain rule for mutual information gives
\[
I_{q,P_1\triangleright\mathcal Q}(A;X_1,X_2)
=
I_{q,P_1}(A;X_1)
+
\sum_{x:m(x)>0}m(x)\,I_{r_x,Q^x}(A;X_2^x).
\tag{N3}
\]
Therefore, if
\[
r_x^0\succeq^I_{Q^x\sqcup R^x}r_x^1
\quad\text{for every positive-probability branch }x,
\]
then $I_{r_x,Q^x}\ge I_{r_x,R^x}$ for every such $x$.  Using (N3),
\[
I_{q,P_1\triangleright\{Q^x\}}
\ge
I_{q,P_1\triangleright\{R^x\}},
\]
so
\[
q^0
\succeq^I_{(P_1\triangleright\{Q^x\})\sqcup(P_1\triangleright\{R^x\})}
q^1.
\]
If one branch comparison is strict and that branch has positive probability, then the corresponding summand in (N3) is
strictly larger, so the aggregate comparison is strict.

No product condition is imposed, so the verification ends here.  For later use, write $q_1\otimes q_2$ for the product
lottery and define the product of channels by
\[
(P_1\otimes P_2)(x_1,x_2\mid a_1,a_2)
:=P_1(x_1\mid a_1)P_2(x_2\mid a_2).
\]
The mutual-information family satisfies the background-inertness property derived in Lemma~\ref{pt:lem:background} below,
since
\[
I_{q_1\otimes q_2,P_1\otimes R_2}(A_1,A_2;X_1,X_2)
=
I_{q_1,P_1}(A_1;X_1)+I_{q_2,R_2}(A_2;X_2),
\tag{N4}
\]
so that the comparison of $P_1\otimes R_2$ and $Q_1\otimes R_2$ at $q_1\otimes q_2$ reduces to the comparison of $P_1$
and $Q_1$ at $q_1$, whatever the common background $R_2$.  Consistently with Lemma~\ref{pt:lem:background}, this is a
consequence of the chain rule (N3) and not an independent requirement.

\subsubsection{Sufficiency}

Assume Conditions~\textup{(P1)--(P5)}.  We prove the implication \textup{(i)$\Rightarrow$(ii)} in
Proposition~\ref{pt:prop:main}.  The proof is written first for full-support lotteries.  Boundary cases are obtained by
exact restriction to the positive support and transport back to the original action set.  The key point is that posterior-law sufficiency is not assumed; it
is derived from block comparisons and record post-processing, and so is its continuity
(Lemma~\ref{pt:lem:plcont}).

\paragraph{How to read the proof.}
The argument has five interfaces.  Each stage ends with exactly the object
needed by the next one:
\begin{enumerate}[label=\textup{Stage \arabic*:},leftmargin=*]
\item replace channels by Bayes-plausible posterior laws and derive the needed
  quotient continuity from primitive closedness;
\item cardinalise each fixed-prior quotient and make support, neutrality, and
  relabelling behaviour explicit;
\item transport those affine values through independent products and into one
  cross-prior block scale, allowing provisionally a bilinear interaction;
\item identify reached-branch coefficients, prove their cocycle, and align
  the scales on nested support faces;
\item compare product and sequential decompositions, eliminate the
  interaction, and invoke Faddeev's recursion to obtain Shannon entropy and
  hence mutual information.
\end{enumerate}
The long posterior-grid, product-coefficient, and branch tangent-space
calculations are technical subarguments.  A reader interested in the logical
spine may use their lemma statements as interfaces.  The kernel proof chooses
its canonical representative before the branch and product calculations and
implements the branch interfaces using finite atomic signed posterior laws.
The paper keeps product calibration before branch calibration because this
makes the scalar algebra easier to read.  The early full-revelation
normalisation and exact relabelling lemma below supply the coherence needed by
that presentation; the kernel proof reaches the same later interfaces in the
opposite order.

\paragraph{Stage 1: posterior-law quotient and derived continuity.}
\addcontentsline{toc}{subsubsection}{Stage 1: posterior-law quotient and derived continuity}
The first five lemmas establish a coherent order on posterior laws.  In
particular, posterior-law sufficiency and continuity are conclusions, not
extra axioms; the barycentric spread--merge construction in
Lemma~\ref{pt:lem:plcont} is the technical core of this stage.

\begin{lemma}[Coherent block comparisons]\label{pt:lem:blockcoh}
Assume Conditions~\textup{(P1)} and \textup{(P3)}.

\textup{(a)}  Let $(q_1,P_1)$, $(q_2,P_2)$, $(q_3,P_3)$ be pairs of a lottery $q_k\in\Delta(A_k)$ and a channel
$P_k:A_k\to\Delta(X_k)$, with action sets, record sets, and lotteries not necessarily equal and the lotteries not
necessarily of full support.  Then the two-block comparisons between such pairs are complete,
\[
q_1^0\succeq_{P_1\sqcup P_2}q_2^1
\quad\text{or}\quad
q_2^0\succeq_{P_2\sqcup P_1}q_1^1,
\]
satisfy the reverse-block equivalence
\[
q_2^1\succeq_{P_1\sqcup P_2}q_1^0
\quad\Longleftrightarrow\quad
q_2^0\succeq_{P_2\sqcup P_1}q_1^1,
\]
and are transitive:
\[
q_1^0\succeq_{P_1\sqcup P_2}q_2^1
\quad\text{and}\quad
q_2^0\succeq_{P_2\sqcup P_3}q_3^1
\quad\Longrightarrow\quad
q_1^0\succeq_{P_1\sqcup P_3}q_3^1.
\]
In particular, for a fixed full-support $q\in\Delta(A)$ the comparisons $q^0\succeq_{P\sqcup Q}q^1$ define a weak order
over channels on $A$ at prior $q$, and pairwise comparisons between block-supported lotteries are independent of the
block environment in which the two blocks are embedded.

\textup{(b)}  Under \textup{(P2)}, the comparisons in \textup{(a)} are also closed in the channels at fixed alphabets
and fixed lotteries: if $P_n\to P$ in $\Delta(X_1)^{A_1}$, $Q_n\to Q$ in $\Delta(X_2)^{A_2}$, $q_1\in\Delta(A_1)$ and
$q_2\in\Delta(A_2)$ are fixed, and
\[
q_1^0\succeq_{P_n\sqcup Q_n}q_2^1
\quad\text{for all }n,
\]
then
\[
q_1^0\succeq_{P\sqcup Q}q_2^1.
\]
Fixedness of the alphabets is inherited from Condition~\textup{(P2)}, which is stated at a fixed pair of alphabets.
\end{lemma}

\begin{proof}
\textup{(a)}  Completeness follows from Condition~\textup{(P1)} applied to the two-block environment $P_1\sqcup P_2$: either
$q_1^0\succeq_{P_1\sqcup P_2}q_2^1$ or $q_2^1\succeq_{P_1\sqcup P_2}q_1^0$.  In the second case, applying
Condition~\textup{(P3)} to the two-block environment with blocks $P_1$ and $P_2$, using the ordered pair $(i,j)=(1,0)$,
gives the canonical reverse comparison
\[
q_2^1\succeq_{P_1\sqcup P_2}q_1^0
\quad\Longleftrightarrow\quad
q_2^0\succeq_{P_2\sqcup P_1}q_1^1,
\]
which is the reverse-block equivalence.

For transitivity, form the three-block environment $\bigsqcup_{k=1}^3 P_k$.  By the irrelevant-blocks clause of
Condition~\textup{(P3)} with block pair $(1,2)$, the first hypothesis is equivalent to
\[
q_1^1\succeq_{\bigsqcup_{k=1}^3 P_k}q_2^2,
\]
and by the same clause with block pair $(2,3)$, the second hypothesis is equivalent to
\[
q_2^2\succeq_{\bigsqcup_{k=1}^3 P_k}q_3^3.
\]
Transitivity of the weak order given by Condition~\textup{(P1)} in this common environment gives
\[
q_1^1\succeq_{\bigsqcup_{k=1}^3 P_k}q_3^3,
\]
and applying Condition~\textup{(P3)} again, now deleting the irrelevant middle block, yields
$q_1^0\succeq_{P_1\sqcup P_3}q_3^1$.  Nothing in these arguments requires the pairs to share alphabets or lotteries, or
the lotteries to have full support.  The specialisation to a fixed full-support $q$ and channels on a common $A$ is
immediate.

\textup{(b)}  The block channel $P_n\sqcup Q_n$ has fixed tagged copies of the alphabets $A_1,A_2,X_1,X_2$ and converges
to $P\sqcup Q$ in the finite-dimensional product simplex.  Since
the block-supported lotteries $q_1^0,q_2^1$ are fixed, Condition~\textup{(P2)} applied at this pair of alphabets gives
\[
q_1^0\succeq_{P\sqcup Q}q_2^1.
\qedhere
\]
\end{proof}

\begin{lemma}[Prior-specific Blackwell equivalence]\label{pt:lem:blackwell}
Let $q\in\Delta(A)$ have full support.  Let $P:A\to\Delta(X)$ and $Q:A\to\Delta(Y)$ be finite channels.  If
\[
\mu_{q,P}=\mu_{q,Q},
\]
then there are stochastic kernels $T:X\to\Delta(Y)$ and $T':Y\to\Delta(X)$ such that
\[
Q=PT,
\qquad
P=QT'.
\]
\end{lemma}

\begin{proof}
It is enough to construct $T$; the construction of $T'$ is symmetric.  Let
\[
m_P(x):=m_{q,P}(x),\qquad m_Q(y):=m_{q,Q}(y).
\]
If $m_P(x)=0$, then $P(x\mid a)=0$ for every $a\in A$, because $q$ has full support.  Such records do not affect
$PT$, so define the corresponding row of $T$ arbitrarily.  The same observation applies to null records of $Q$.

Let $\mathcal R$ be the finite set of posterior vectors that occur with positive probability under either channel.  For
$r\in\mathcal R$, define
\[
X_r:=\{x\in X:m_P(x)>0,\ r_x^{\,q,P}=r\},
\qquad
Y_r:=\{y\in Y:m_Q(y)>0,\ r_y^{\,q,Q}=r\}.
\]
The sets $\{X_r\}_{r\in\mathcal R}$ partition the positive-probability records of $P$, and the sets
$\{Y_r\}_{r\in\mathcal R}$ partition the positive-probability records of $Q$, by posterior vector.
Equality of posterior laws implies that the total probability attached to posterior $r$ is the same for the two
channels:
\[
\lambda_r:=\sum_{x\in X_r}m_P(x)=\sum_{y\in Y_r}m_Q(y)>0.
\]
For $x\in X_r$ and $y\in Y_r$, set
\[
T(y\mid x):=\frac{m_Q(y)}{\lambda_r},
\]
and set $T(y\mid x)=0$ when $x\in X_r$ and $y\notin Y_r$.  This defines a stochastic row for every positive-probability
record $x$, since the probabilities over $Y_r$ sum to one.  Rows corresponding to null records of $P$ were already
chosen arbitrarily.

Fix $a\in A$ and $y\in Y_r$.  Since all records in $X_r$ have posterior $r$, Bayes' rule gives
\[
q(a)P(x\mid a)=m_P(x)r(a)
\qquad(x\in X_r).
\]
Therefore
\[
\sum_{x\in X_r}P(x\mid a)
=
\frac{\lambda_r r(a)}{q(a)}.
\]
Using the definition of $T$,
\[
(PT)(y\mid a)
=
\sum_{x\in X_r}P(x\mid a)\frac{m_Q(y)}{\lambda_r}
=
\frac{m_Q(y)r(a)}{q(a)}.
\]
But $y\in Y_r$ also has posterior $r$, so Bayes' rule for $Q$ gives
\[
Q(y\mid a)=\frac{m_Q(y)r(a)}{q(a)}.
\]
Thus $(PT)(y\mid a)=Q(y\mid a)$ for every positive-probability $y$.  If $m_Q(y)=0$, then full support of $q$ implies
$Q(y\mid a)=0$ for every $a$.  No positive-probability record of $P$ sends mass to such a $y$ under the construction;
rows of $T$ indexed by null records of $P$ may be chosen arbitrarily, but those records have $P(x\mid a)=0$ for every
$a$, so they contribute nothing to $PT$.  Hence $(PT)(y\mid a)=0=Q(y\mid a)$ for every $a$, and therefore $Q=PT$.
\end{proof}

\begin{lemma}[Posterior-law sufficiency]\label{pt:lem:plsuff}
Fix a full-support $q\in\Delta(A)$.  If $P:A\to\Delta(X)$ and $Q:A\to\Delta(Y)$ generate the same posterior law at $q$,
\[
\mu_{q,P}=\mu_{q,Q},
\]
then
\[
q^0\sim_{P\sqcup Q}q^1.
\]
Consequently, by Lemma~\ref{pt:lem:blockcoh}, at a fixed full-support prior the block comparison of channels depends only
on the induced posterior law.
\end{lemma}

\begin{proof}
By Lemma~\ref{pt:lem:blackwell}, there are stochastic kernels $T$ and $T'$ such that $Q=PT$ and $P=QT'$.  Clause
\textup{(DP-rec)} of Condition~\textup{(P4)} gives
\[
q^0\succeq_{P\sqcup PT}q^1,
\qquad
q^0\succeq_{Q\sqcup QT'}q^1.
\]
Using $PT=Q$, the first comparison is
\[
q^0\succeq_{P\sqcup Q}q^1.
\]
Using $QT'=P$, the second comparison is
\[
q^0\succeq_{Q\sqcup P}q^1.
\]
By the reverse-block equivalence in Lemma~\ref{pt:lem:blockcoh} (equivalently, Condition~\textup{(P3)} applied to
$P\sqcup Q$ with ordered pair $(i,j)=(1,0)$), this is equivalent to
\[
q^1\succeq_{P\sqcup Q}q^0.
\]
Thus $q^0\sim_{P\sqcup Q}q^1$.
\end{proof}

For channels $P:A\to\Delta(X_P)$ and $R:A\to\Delta(X_R)$ and $\lambda\in[0,1]$, write
$\lambda P\oplus(1-\lambda)R$ for the channel that publicly selects $P$ with probability $\lambda$ and $R$ otherwise,
and reports the selected channel together with its record.  Thus
\[
\bigl(\lambda P\oplus(1-\lambda)R\bigr)((z,i)\mid a)
=
\begin{cases}
\lambda P(z\mid a),&i=0,\\
(1-\lambda)R(z\mid a),&i=1.
\end{cases}
\]
Because the coin is observed,
\[
\mu_{q,\lambda P\oplus(1-\lambda)R}
=\lambda\mu_{q,P}+(1-\lambda)\mu_{q,R}.
\]

\begin{lemma}[Derived public-coin independence]\label{pt:lem:publiccoin}
Assume Conditions~\textup{(P1)} and \textup{(P3)}, clause \textup{(DP-rec)} of Condition~\textup{(P4)}, and
Condition~\textup{(P5)}.  Fix a finite action set $A$, a
full-support $q\in\Delta(A)$, and channels $P,Q,R$ on $A$.  For every $\lambda\in(0,1)$,
\[
q^0\succeq_{P\sqcup Q}q^1
\quad\Longleftrightarrow\quad
q^0\succeq_{(\lambda P\oplus(1-\lambda)R)\sqcup(\lambda Q\oplus(1-\lambda)R)}q^1.
\]
\end{lemma}

\begin{proof}
Throughout, for channels $V,W$ on $A$ write $V\succeq_q W$ for the block comparison $q^0\succeq_{V\sqcup W}q^1$, and
$V\succ_q W$ for the corresponding strict comparison.  By Lemma~\ref{pt:lem:blockcoh} the comparison $\succeq_q$ is
complete and transitive, and by Lemma~\ref{pt:lem:plsuff} it depends only on the pair of posterior laws
$(\mu_{q,V},\mu_{q,W})$.  In particular, channels with equal posterior laws at $q$ may be substituted on either side of
any comparison; we refer to this as \emph{substitution}.  Substitution also preserves strict comparisons: if
$\mu_{q,V}=\mu_{q,V'}$ and $\mu_{q,W}=\mu_{q,W'}$, then $V\sim_q V'$ and $W\sim_q W'$ by Lemma~\ref{pt:lem:plsuff}, and
transitivity of the complete relation $\succeq_q$ gives $V\succ_q W$ if and only if $V'\succ_q W'$.

\emph{Step 1: the public coin is an uninformative first stage.}
Define the constant channel $C:A\to\Delta(\{0,1\})$ by
\[
C(0\mid a):=\lambda,
\qquad
C(1\mid a):=1-\lambda
\qquad(a\in A).
\]
Since $\lambda\in(0,1)$, both branches have positive mass: $m_{q,C}(0)=\lambda$ and $m_{q,C}(1)=1-\lambda$.  Because
the rows of $C$ are constant, Bayes' rule gives, for each $x\in\{0,1\}$,
\[
r_x^{\,q,C}(a)=\frac{q(a)C(x\mid a)}{m_{q,C}(x)}=q(a),
\]
so both branch posteriors equal the prior: $r_0=r_1=q$.

\emph{Step 2: padding the continuations to a common alphabet.}
Condition~\textup{(P5)} requires the two continuation channels on a given branch to share a record alphabet, whereas $P$
and $Q$ may not.  Let $X_{PQ}:=X_P\sqcup X_Q$ and define $\widetilde P,\widetilde Q:A\to\Delta(X_{PQ})$ by
\[
\widetilde P(z\mid a):=
\begin{cases}
P(z\mid a)&z\in X_P,\\
0&z\in X_Q,
\end{cases}
\qquad
\widetilde Q(z\mid a):=
\begin{cases}
0&z\in X_P,\\
Q(z\mid a)&z\in X_Q.
\end{cases}
\]
Padding adds only zero-probability records, so the positive-mass records and their posteriors are unchanged:
$\mu_{q,\widetilde P}=\mu_{q,P}$ and $\mu_{q,\widetilde Q}=\mu_{q,Q}$.  By substitution,
\[
q^0\succeq_{P\sqcup Q}q^1
\quad\Longleftrightarrow\quad
\widetilde P\succeq_q\widetilde Q,
\]
and the same equivalence holds for the strict comparisons.

\emph{Step 3: the compounds implement the public-coin mixtures.}
Consider the two-stage channels with first stage $C$, continuation pair $\widetilde P$ versus $\widetilde Q$ on branch
$0$, and the common continuation $R$ on branch $1$:
\[
C\triangleright\{\widetilde P,R\}
\qquad\text{and}\qquad
C\triangleright\{\widetilde Q,R\},
\]
where the profile notation lists the branch-$0$ and branch-$1$ continuations in order.  By the definition of two-stage
composition,
\[
(C\triangleright\{\widetilde P,R\})(0,z\mid a)=\lambda\,\widetilde P(z\mid a),
\qquad
(C\triangleright\{\widetilde P,R\})(1,z\mid a)=(1-\lambda)\,R(z\mid a).
\]
A record $(0,z)$ has marginal mass $\lambda\,m_{q,\widetilde P}(z)$, and when this is positive the constant factor
$\lambda$ cancels in Bayes' rule, so its posterior is $r_z^{\,q,\widetilde P}$; symmetrically for records $(1,z)$.
Hence
\[
\mu_{q,C\triangleright\{\widetilde P,R\}}
=
\lambda\mu_{q,\widetilde P}+(1-\lambda)\mu_{q,R}
=
\lambda\mu_{q,P}+(1-\lambda)\mu_{q,R}
=
\mu_{q,\lambda P\oplus(1-\lambda)R},
\]
using the public-coin mixing identity, and likewise
$\mu_{q,C\triangleright\{\widetilde Q,R\}}=\mu_{q,\lambda Q\oplus(1-\lambda)R}$.  By substitution, the right-hand
comparison of the lemma is therefore equivalent to
\[
C\triangleright\{\widetilde P,R\}
\succeq_q
C\triangleright\{\widetilde Q,R\}.
\tag{$\ast$}
\]

\emph{Step 4: forward implication.}
Assume $q^0\succeq_{P\sqcup Q}q^1$.  We verify the branch hypotheses of Condition~\textup{(P5)} for the first stage $C$ and
the continuation profiles $\{\widetilde P,R\}$ versus $\{\widetilde Q,R\}$.  Both branches have positive mass and both
branch posteriors equal $q$ (Step~1).  On branch $0$ the hypothesis is
$q^0\succeq_{\widetilde P\sqcup\widetilde Q}q^1$, which holds by Step~2.  On branch $1$ the two profiles share the
continuation $R$, and the hypothesis $q^0\succeq_{R\sqcup R}q^1$ holds because Condition~\textup{(P1)} gives
$q\succeq_R q$ and the duplication clause of Condition~\textup{(P3)} converts it into the two-block form.
Condition~\textup{(P5)} then yields $(\ast)$, hence the mixture comparison by Step~3.

\emph{Step 5: converse, via the strictness clause.}
Assume the mixture comparison holds, so that $(\ast)$ holds by Step~3, and suppose towards a contradiction that
$q^0\succeq_{P\sqcup Q}q^1$ fails.  By completeness and the reverse-block equivalence of Lemma~\ref{pt:lem:blockcoh},
$Q\succ_q P$, and Step~2 gives $\widetilde Q\succ_q\widetilde P$.  Now apply Condition~\textup{(P5)} with the roles of the
two profiles exchanged: the branch-$0$ comparison $\widetilde Q\succ_q\widetilde P$ is strict on a branch of mass
$\lambda>0$, and the branch-$1$ comparison $R\succeq_q R$ holds as in Step~4.  The strictness clause of
Condition~\textup{(P5)} yields
\[
C\triangleright\{\widetilde Q,R\}
\succ_q
C\triangleright\{\widetilde P,R\},
\]
which contradicts $(\ast)$.  Hence $q^0\succeq_{P\sqcup Q}q^1$.
\end{proof}

The lemma shows that public-coin independence is not an independent behavioural commitment: a public coin is exactly a
first-stage observation that reveals nothing about the act, so invariance under common public mixtures is branchwise
continuation monotonicity applied through trivial branch posteriors.  The statement is given at full-support priors,
which is the only case used below (Lemma~\ref{pt:lem:postsep}).

For a finite action set $A$, let $\mathsf X_A$ be the set of finite-support probability measures on $\Delta(A)$.  Write
\[
b(\mu):=\int_{\Delta(A)}r\,d\mu(r),
\qquad
\mathcal M_q:=\{\mu\in\mathsf X_A:b(\mu)=q\}.
\]
Thus $\mathcal M_q$ is the set of finite posterior laws with prior $q$ as their mean, and
$\mu_{q,P}\in\mathcal M_q$ for every channel $P$.

\begin{lemma}[Posterior-law continuity]\label{pt:lem:plcont}
Fix a full-support $q\in\Delta(A)$.  Let $P_n,Q_n,P,Q$ be finite channels on $A$ with
\[
\mu_{q,P_n}\Rightarrow \mu_{q,P},
\qquad
\mu_{q,Q_n}\Rightarrow \mu_{q,Q}.
\]
If
\[
q^0\succeq_{P_n\sqcup Q_n}q^1
\qquad\text{for all }n,
\]
then
\[
q^0\succeq_{P\sqcup Q}q^1.
\]
\end{lemma}

\begin{proof}
Throughout the proof, for channels $V,W$ on $A$ write $V\succeq_q W$ for the block comparison
$q^0\succeq_{V\sqcup W}q^1$.  By Lemma~\ref{pt:lem:blockcoh} this comparison is complete and transitive, and by
Lemma~\ref{pt:lem:plsuff} it depends only on the pair of posterior laws $(\mu_{q,V},\mu_{q,W})$; in particular, channels
with equal posterior laws at $q$ may be substituted on either side of any comparison.  We refer to this as
\emph{substitution}.

For a finite-support Bayes-plausible law
\[
\mu=\sum_{k\in K}\lambda_k\delta_{r_k}\in\mathcal M_q,
\]
with atoms listed so that $\lambda_k>0$, define the canonical implementing channel
$C_\mu:A\to\Delta(K)$ by
\[
C_\mu(k\mid a):=\frac{\lambda_k r_k(a)}{q(a)}.
\]
Full support of $q$ and $\sum_k\lambda_k r_k=q$ make $C_\mu$ a channel with $\mu_{q,C_\mu}=\mu$.

\emph{Step 1: merging posteriors moves weakly down.}
Let $\mu=\sum_{k\in K}\lambda_k\delta_{r_k}\in\mathcal M_q$ and let $\{K_1,\dots,K_m\}$ partition $K$ into nonempty
groups.  The merged law is
\[
\operatorname{mrg}(\mu):=\sum_{j=1}^m w_j\,\delta_{s_j},
\qquad
w_j:=\sum_{k\in K_j}\lambda_k,
\qquad
s_j:=\frac{1}{w_j}\sum_{k\in K_j}\lambda_k r_k.
\]
Barycentres are preserved, so $\operatorname{mrg}(\mu)\in\mathcal M_q$.  Let $T:K\to\Delta(\{1,\dots,m\})$ be the
deterministic kernel with $T(j\mid k)=1$ if and only if $k\in K_j$.  Then
\[
(C_\mu T)(j\mid a)
=
\sum_{k\in K_j}\frac{\lambda_k r_k(a)}{q(a)}
=
\frac{w_j s_j(a)}{q(a)}
=
C_{\operatorname{mrg}(\mu)}(j\mid a),
\]
so clause \textup{(DP-rec)} of Condition~\textup{(P4)} gives
\[
C_\mu\succeq_q C_{\operatorname{mrg}(\mu)}.
\]
Conversely, say that $\mu'\in\mathcal M_q$ is a \emph{spread} of $\mu$ if $\mu'$ is obtained by replacing each atom
$(\lambda_k,r_k)$ of $\mu$ by finitely many atoms $(\lambda_k w_{kj},v_{kj})_j$ with $w_{kj}\ge0$,
$\sum_j w_{kj}=1$, and $\sum_j w_{kj}v_{kj}=r_k$.  Then $\mu=\operatorname{mrg}(\mu')$ for the grouping by $k$, so
\[
C_{\mu'}\succeq_q C_\mu.
\]
Atoms with $w_{kj}=0$ are deleted from the list; and if the same posterior appears in several list entries, any two
listings of the same measure induce the same posterior law at $q$, so the corresponding canonical channels are
interchangeable by substitution.
Merging records can only obscure which act was realised; splitting an atom into a barycentric family is the reverse
operation.

\emph{Step 2: bounded supports.}
We claim the lemma holds under the additional hypothesis that all laws $\mu_{q,P_n}$ and $\mu_{q,Q_n}$ have support of
size at most $K$, for some fixed $K$.  Write
\[
\mu_{q,P_n}=\sum_{k=1}^{K}\lambda_{n,k}\,\delta_{r_{n,k}},
\]
padding with zero-mass atoms ($\lambda_{n,k}=0$, $r_{n,k}:=q$) so that exactly $K$ terms appear, and define the
$K$-record channel $C_n:A\to\Delta(\{1,\dots,K\})$ by
\[
C_n(k\mid a):=\frac{\lambda_{n,k}\,r_{n,k}(a)}{q(a)}.
\]
Each row sums to one because $\sum_k\lambda_{n,k}r_{n,k}=q$, zero-mass columns contribute nothing, and
$\mu_{q,C_n}=\mu_{q,P_n}$.  Define $D_n$ from $\mu_{q,Q_n}$ in the same way.  By substitution,
\[
C_n\succeq_q D_n
\qquad\text{for all }n.
\]
The parameter sequences of $C_n$ and $D_n$ jointly live in the compact metric space
$([0,1]\times\Delta(A))^{2K}$; pass to a single subsequence along which $\lambda_{n,k}\to\lambda_k^\ast$ and
$r_{n,k}\to r_k^\ast$ for every $k$, simultaneously for both channels.  Then $C_n\to C^\ast$ entrywise, where
$C^\ast(k\mid a)=\lambda_k^\ast r_k^\ast(a)/q(a)$; each row of $C^\ast$ sums to one as an entrywise limit of rows
summing to one, so $C^\ast$ is a channel.  For every continuous $f$ on $\Delta(A)$,
\[
\int f\,d\mu_{q,P_n}
=
\sum_{k=1}^{K}\lambda_{n,k}f(r_{n,k})
\longrightarrow
\sum_{k=1}^{K}\lambda_k^\ast f(r_k^\ast),
\]
so by uniqueness of weak limits
\[
\mu_{q,P}=\sum_{k=1}^{K}\lambda_k^\ast\,\delta_{r_k^\ast}
=\mu_{q,C^\ast},
\]
where the second equality holds because record $k$ of $C^\ast$ has probability $\lambda_k^\ast$ and, when
$\lambda_k^\ast>0$, posterior $r_k^\ast$.  Similarly $D_n\to D^\ast$ entrywise with
$\mu_{q,D^\ast}=\mu_{q,Q}$.  The closedness clause, Lemma~\ref{pt:lem:blockcoh}(b), applied along the subsequence gives
\[
C^\ast\succeq_q D^\ast,
\]
and substitution returns $q^0\succeq_{P\sqcup Q}q^1$.  This proves the bounded-support case; note that only the
primitive closedness of Condition~\textup{(P2)} has been used, through Lemma~\ref{pt:lem:blockcoh}.

\emph{Step 3: the general case, by an $\varepsilon$-grid sandwich.}
Write $\mu_n:=\mu_{q,P_n}$, $\nu_n:=\mu_{q,Q_n}$, $\mu:=\mu_{q,P}$, $\nu:=\mu_{q,Q}$.  The only remaining gap is that
the supports of $\mu_n,\nu_n$ need not be bounded.  Fix $\varepsilon>0$.  Two finite partitions of $\Delta(A)$ are
used, one for each surrogate; they play independent roles and need not be related.  For the spread, fix a finite
triangulation $\mathcal T$ of $\Delta(A)$ into simplices of diameter at most $\varepsilon$, together with a measurable
partition $\{\widetilde S:S\in\mathcal T\}$ of $\Delta(A)$ with $\widetilde S\subseteq S$, assigning each shared face
to exactly one incident cell.  For the merge, fix a finite Borel partition $\{E_1,\dots,E_M\}$ of $\Delta(A)$ into
sets of diameter at most $\varepsilon$ whose relative boundaries carry no $\nu$-mass: cover $\Delta(A)$ by finitely
many open balls of radius at most $\varepsilon/2$, choosing the radii so that no atom of $\nu$ lies on a bounding
sphere ($\nu$ has finitely many atoms, so all but finitely many radii work), and disjointify.

For finite-support $\rho\in\mathcal M_q$ define two bounded-support surrogates.  Both remain in
$\mathcal M_q$: the spread operation replaces each posterior by a barycentric decomposition of itself, while the merge
operation replaces a group of posteriors by their conditional mean, weighted by the total mass of the group.  Cells with
zero mass are ignored (or, equivalently, padded by an arbitrary posterior with zero weight when a fixed record alphabet is
convenient).
The \emph{grid spread} $\operatorname{sp}_\varepsilon(\rho)$ replaces each atom $(\lambda,\beta)$ with
$\beta\in\widetilde S$ by the atoms $\bigl(\lambda\,w_v(\beta),\,v\bigr)_{v}$, where $v$ ranges over the vertices of
$S$ and $(w_v(\beta))_v$ are the barycentric coordinates of $\beta$ in $S$.  It is a spread of $\rho$ in the sense of
Step~1, is supported on the fixed vertex set of $\mathcal T$, and satisfies
$C_{\operatorname{sp}_\varepsilon(\rho)}\succeq_q C_\rho$.
The \emph{cell merge} $\operatorname{mg}_\varepsilon(\rho)$ merges the atoms of $\rho$ within each class $E_j$ of the
Borel partition; it has at most one atom per class and satisfies
$C_\rho\succeq_q C_{\operatorname{mg}_\varepsilon(\rho)}$ by Step~1, since Step~1 allows any grouping of atoms.

By substitution and transitivity (Lemma~\ref{pt:lem:blockcoh}), the hypothesis $P_n\succeq_q Q_n$ gives the chain
\[
C_{\operatorname{sp}_\varepsilon(\mu_n)}
\succeq_q
C_{\mu_n}
\succeq_q
C_{\nu_n}
\succeq_q
C_{\operatorname{mg}_\varepsilon(\nu_n)}.
\]

Both surrogate sequences converge weakly on fixed alphabets.  For the spread side, let $f$ be continuous on
$\Delta(A)$.  The triangulation defines a continuous piecewise-affine interpolant $I_{\mathcal T}f$ by
\[
I_{\mathcal T}f(r):=\sum_{v\in\operatorname{vert}(S)}w_v(r)f(v)
\qquad\text{for }r\in S.
\]
This is well-defined: on a shared face the barycentric coordinates computed in the two incident simplices put zero
weight on vertices outside the face and coincide with the intrinsic barycentric coordinates on that face.  Hence
$I_{\mathcal T}f$ is continuous on all of $\Delta(A)$.  Moreover, for every finite-support $\rho$,
\[
\int f\,d\operatorname{sp}_\varepsilon(\rho)
= \int I_{\mathcal T}f\,d\rho,
\]
because each posterior is replaced by the barycentric lottery over the vertices of its simplex.  Thus
$\mu_n\Rightarrow\mu$ gives
$\operatorname{sp}_\varepsilon(\mu_n)\Rightarrow\operatorname{sp}_\varepsilon(\mu)$.  For the merge side, each class
$E_j$ has relative boundary contained in the chosen bounding spheres, which carry no $\nu$-mass, so $E_j$ is a
$\nu$-continuity set; therefore
$\nu_n(E_j)\to\nu(E_j)$ and, since the truncated map $r\mapsto r\,\1_{E_j}(r)$ is bounded and its discontinuity set
lies in the relative boundary of $E_j$, which is $\nu$-null,
$\int_{E_j}r\,d\nu_n\to\int_{E_j}r\,d\nu$ for every
class.  Classes with $\nu(E_j)>0$ thus have converging merged masses and merged locations.  If $\nu(E_j)=0$, then the
merged mass tends to zero, so the chosen location in that class is immaterial for weak convergence; padding such a class
by any fixed posterior only changes a zero-weight atom.  Hence
$\operatorname{mg}_\varepsilon(\nu_n)\Rightarrow\operatorname{mg}_\varepsilon(\nu)$.

The laws $\operatorname{sp}_\varepsilon(\mu_n)$ and $\operatorname{mg}_\varepsilon(\nu_n)$ have supports of size at
most $N(\varepsilon)$, the maximum of the number of vertices of $\mathcal T$ and the number $M$ of classes of the
Borel partition.  Step~2 applied
to the chain's endpoints therefore yields
\[
C_{\operatorname{sp}_\varepsilon(\mu)}
\succeq_q
C_{\operatorname{mg}_\varepsilon(\nu)}.
\]

Finally, send $\varepsilon\to0$ along a sequence $\varepsilon_m\downarrow0$ with triangulations and Borel partitions
chosen as above.  The
law $\operatorname{sp}_{\varepsilon_m}(\mu)$ has at most $k_P\,|A|$ atoms, where $k_P$ is the number of atoms of
$\mu$: each atom splits among the $|A|$ vertices of one cell.  Each new atom lies within $\varepsilon_m$ of the atom
it came from and group masses are preserved, so $\operatorname{sp}_{\varepsilon_m}(\mu)\Rightarrow\mu$.  The law
$\operatorname{mg}_{\varepsilon_m}(\nu)$ has at most $k_Q$ atoms, where $k_Q$ is the number of atoms of $\nu$: only
classes meeting the support of $\nu$ carry mass.  Once $\varepsilon_m$ is smaller than the minimal distance between
distinct atoms of $\nu$, each class contains at most one atom of $\nu$ and
$\operatorname{mg}_{\varepsilon_m}(\nu)=\nu$ exactly.  The comparisons
$C_{\operatorname{sp}_{\varepsilon_m}(\mu)}\succeq_q C_{\operatorname{mg}_{\varepsilon_m}(\nu)}$ hold for every $m$,
the supports are uniformly bounded, and the laws converge weakly to $\mu$ and $\nu$; one further application of
Step~2 gives $C_\mu\succeq_q C_\nu$.  Substitution returns
\[
q^0\succeq_{P\sqcup Q}q^1.
\qedhere
\]
\end{proof}

The lemma shows that continuity of block comparisons in posterior-law convergence is not an independent behavioural
commitment: it is the primitive closedness of Condition~\textup{(P2)} seen through the quotient that
Condition~\textup{(P3)} and clause \textup{(DP-rec)} of Condition~\textup{(P4)} impose, with the convex order supplying compactness where record alphabets
grow without bound.  Only weak preference is closed in the limit; no strictness claim is made or used.

\paragraph{Stage 2: affine fibres and canonical boundary behaviour.}
\addcontentsline{toc}{subsubsection}{Stage 2: affine fibres and canonical boundary behaviour}
Public-coin independence and quotient continuity now permit
Herstein--Milnor cardinalisation.  The next lemmas record the finite integral
form, normalize no information to zero, and make action coarsening, exact
support restriction, and undetectable distinctions available before values
are compared across priors.  In the formal development these facts are
packaged by a full-revelation-normalized selected representative.  The paper
now makes that same initial selection explicitly; a later coherent product
gauge is declared when it changes the representative.

\begin{lemma}[Posterior-separable representation]\label{pt:lem:postsep}
Fix a full-support $q\in\Delta(A)$.  There is a continuous affine functional
\[
F_q:\mathcal M_q\to\R
\]
such that, for any two channels $P,Q$ on $A$,
\[
q^0\succeq_{P\sqcup Q}q^1
\quad\Longleftrightarrow\quad
F_q(\mu_{q,P})\ge F_q(\mu_{q,Q}).
\]
Moreover, for some continuous $\phi_q:\Delta(A)\to\R$,
\[
F_q(\mu)=\int_{\Delta(A)}\phi_q(r)\,d\mu(r),
\]
and therefore
\[
F_q(\mu_{q,P})
=
\sum_{x:m_{q,P}(x)>0}m_{q,P}(x)\,\phi_q(r_x^{\,q,P}).
\]
The affine functional is unique up to positive affine transformations.  The integrand $\phi_q$ is not pointwise unique;
adding an affine function of $r$ leaves its integral on the fibre $\mathcal M_q$ changed only by a constant.
\end{lemma}

\begin{proof}
First note that every finite-support Bayes-plausible law is implemented by a finite channel.  If
\[
\mu=\sum_{k=1}^n\lambda_k\delta_{r_k}\in\mathcal M_q,
\]
define
\[
P(k\mid a)=\frac{\lambda_k r_k(a)}{q(a)}.
\]
Since $q$ has full support and $\sum_k\lambda_k r_k=q$, this is a channel.  Its record probability at $k$ is
$\lambda_k$, and the posterior after $k$ is $r_k$; hence $\mu_{q,P}=\mu$.

Define a comparison on $\mathcal M_q$ as follows: for $\mu,\nu\in\mathcal M_q$, choose any channels $P,Q$ implementing
them and compare the two block-supported copies of $q$ in $P\sqcup Q$.  This comparison is independent of the chosen
implementations.  Indeed, if $P,P'$ both implement $\mu$ and $Q,Q'$ both implement $\nu$, then
Lemma~\ref{pt:lem:plsuff} gives
\[
q^0\sim_{P\sqcup P'}q^1,
\qquad
q^0\sim_{Q\sqcup Q'}q^1.
\]
Placing the relevant blocks in common block environments and using Lemma~\ref{pt:lem:blockcoh}, transitivity gives
\[
q^0\succeq_{P\sqcup Q}q^1
\quad\Longleftrightarrow\quad
q^0\succeq_{P'\sqcup Q'}q^1.
\]
Thus the quotient comparison on $\mathcal M_q$ is well defined; denote it by $\succeq_q$.  This construction is only at
the fixed full-support prior $q$, and every law in $\mathcal M_q$ is finite-support and, by the preceding construction,
implemented by a finite channel.  Lemma~\ref{pt:lem:blockcoh} also gives completeness and transitivity of this quotient
comparison.

The set $\mathcal M_q$ is a mixture space under ordinary convex combination.  It is convex because barycentres are
preserved by convex combination.  The quotient relation $\succeq_q$ is a weak order by Lemma~\ref{pt:lem:blockcoh}.
Public-coin mixtures of implementing channels correspond exactly to convex mixtures of posterior laws:
\[
\mu_{q,\lambda P\oplus(1-\lambda)R}
=
\lambda\mu_{q,P}+(1-\lambda)\mu_{q,R}.
\]
Therefore the derived public-coin independence property (Lemma~\ref{pt:lem:publiccoin}) gives mixture independence: for every $\mu,\nu,\rho\in\mathcal M_q$ and
$\lambda\in(0,1)$,
\[
\mu\succeq_q\nu
\quad\Longleftrightarrow\quad
\lambda\mu+(1-\lambda)\rho
\succeq_q
\lambda\nu+(1-\lambda)\rho.
\]

The relative weak topology on $\mathcal M_q$ is inherited from the weak topology on probability measures over the compact
metric simplex $\Delta(A)$, and convex-mixture paths are continuous in this topology.  Lemma~\ref{pt:lem:plcont} gives
relatively closed upper and lower contour sets.  Indeed, if
$\mu_n\Rightarrow\mu$ with $\mu_n,\mu,z\in\mathcal M_q$ and $\mu_n\succeq_q z$, choose finite implementing channels for
these laws and apply Lemma~\ref{pt:lem:plcont} (with the constant sequence implementing $z$) to get $\mu\succeq_q z$; the
argument for $\{\,\mu:z\succeq_q\mu\,\}$ is the
same.  Therefore, for fixed $\mu,\nu,\rho\in\mathcal M_q$, the sets
\[
\{\lambda\in[0,1]:\lambda\mu+(1-\lambda)\nu\succeq_q\rho\}
\quad\text{and}\quad
\{\lambda\in[0,1]:\rho\succeq_q\lambda\mu+(1-\lambda)\nu\}
\]
are closed.

By the Herstein--Milnor mixture-space theorem \cite{HersteinMilnor1953}, there is a mixture-preserving representation
$F_q:\mathcal M_q\to\R$:
\[
\mu\succeq_q\nu
\quad\Longleftrightarrow\quad
F_q(\mu)\ge F_q(\nu),
\]
and
\[
F_q(\lambda\mu+(1-\lambda)\nu)
=
\lambda F_q(\mu)+(1-\lambda)F_q(\nu).
\]
If the quotient order is nonconstant, this representative is unique up to positive affine transformations.  If the
quotient order is constant, take $F_q\equiv0$.

Moreover, $F_q$ is continuous in the relative weak topology.  For any $z\in\mathcal M_q$, the sets
\[
\{\mu:F_q(\mu)\ge F_q(z)\}
=
\{\mu:\mu\succeq_q z\}
\quad\text{and}\quad
\{\mu:F_q(\mu)\le F_q(z)\}
=
\{\mu:z\succeq_q\mu\}
\]
are relatively closed by the preceding Lemma~\ref{pt:lem:plcont} argument.  Since $F_q$ is affine, its range on the convex set
$\mathcal M_q$ is an interval.  Hence every upper and lower level set of $F_q$ is relatively closed: levels outside the
range are trivial, and levels inside the range equal $F_q(z)$ for some $z\in\mathcal M_q$.  Thus $F_q$ is both upper and
lower semicontinuous, and therefore continuous.  For any channels $P,Q$ on $A$,
\[
q^0\succeq_{P\sqcup Q}q^1
\quad\Longleftrightarrow\quad
F_q(\mu_{q,P})\ge F_q(\mu_{q,Q}).
\]

The no-information posterior law at $q$ is $\delta_q$.  Write
\[
\chi_q:=\sum_{a\in A}q(a)\delta_{e_a}
\]
for the full-revelation posterior law, where $e_a$ assigns probability one to action $a$.

It remains to write the affine functional in posterior-separable form, without leaving the finite-support domain.  Let
$X:=\Delta(A)$.  Choose
\[
0<\varepsilon<\min_{a\in A}q(a).
\]
For each $r\in X$, define
\[
s_r:=\frac{q-\varepsilon r}{1-\varepsilon}\in\Delta(A)
\]
and
\[
\eta_r
:=
\varepsilon\delta_r+(1-\varepsilon)\sum_{a\in A}s_r(a)\delta_{e_a}.
\]
Then $\eta_r\in\mathcal M_q$, because its barycentre is
\[
\varepsilon r+(1-\varepsilon)s_r=q.
\]
Choose constants $c_a$ such that
\[
\sum_{a\in A}q(a)c_a=F_q(\chi_q),
\]
for instance $c_a=F_q(\chi_q)$ for every $a$.  Define
\[
\phi_q(r):=
\frac{
F_q(\eta_r)
-
(1-\varepsilon)\sum_{a\in A}s_r(a)c_a
}{\varepsilon}.
\]
The map $r\mapsto\eta_r$ is weakly continuous and $F_q$ is continuous, so $\phi_q$ is continuous.

Now take any finite-support law
\[
\mu=\sum_{k=1}^n\lambda_k\delta_{r_k}\in\mathcal M_q.
\]
Since $\sum_k\lambda_k r_k=q$, we have
\[
\sum_{k=1}^n\lambda_k s_{r_k}=q.
\]
Therefore
\[
\sum_{k=1}^n\lambda_k\eta_{r_k}
=
\varepsilon\mu+(1-\varepsilon)\chi_q.
\]
By affinity of $F_q$,
\[
\sum_{k=1}^n\lambda_k F_q(\eta_{r_k})
=
F_q\!\left(\varepsilon\mu+(1-\varepsilon)\chi_q\right)
=
\varepsilon F_q(\mu)+(1-\varepsilon)F_q(\chi_q).
\]
Using $\sum_k\lambda_k s_{r_k}=q$ and $\sum_a q(a)c_a=F_q(\chi_q)$ gives
\[
\sum_{k=1}^n\lambda_k\phi_q(r_k)=F_q(\mu).
\]
This is exactly
\[
F_q(\mu)=\int_X\phi_q(r)\,d\mu(r)
\]
for every finite-support $\mu\in\mathcal M_q$.  Applying this to $\mu_{q,P}$ gives the stated finite-sum formula.  The
affine representative $F_q$ is unique up to positive affine transformations by Herstein--Milnor.  The integrand is not
pointwise unique on a fixed barycentre fibre: if $\ell$ is affine in $r$, then $\int\ell(r)\,d\mu(r)=\ell(q)$ for every
$\mu\in\mathcal M_q$.
\end{proof}

\begin{lemma}[Refinement monotonicity and canonical normalisation]\label{pt:lem:convnorm}
Fix a full-support $q\in\Delta(A)$ and let $F_q,\phi_q$ be as in Lemma~\ref{pt:lem:postsep}.  If a finite law
$\mu\in\mathcal M_q$ is obtained from another finite law $\nu\in\mathcal M_q$ by replacing each posterior of $\nu$ by a
finite Bayes-plausible continuation experiment, then
\[
F_q(\mu)\ge F_q(\nu).
\]
Consequently, $\phi_q$ may be chosen convex.  After replacing $\phi_q$ by an observationally equivalent representative
and subtracting a constant from $F_q$, one may impose
\[
\phi_q(q)=F_q(\delta_q)=0,
\qquad
\phi_q(r)\ge 0\quad\text{for every }r\in\Delta(A),
\]
and therefore
\[
F_q(\mu)\ge 0\quad\text{for every }\mu\in\mathcal M_q.
\]
If $|A|\ge2$, the representative can and will be selected uniquely by the
additional full-revelation normalisation
\[
F_q(\chi_q)=1.
\]
If $|A|=1$, then $\mathcal M_q=\{\delta_q\}$ and we select $F_q\equiv0$.
\end{lemma}

\begin{proof}
Write
\[
\nu=\sum_{i=1}^n m_i\delta_{r_i},
\qquad
\mu=\sum_{i=1}^n m_i\nu_i,
\qquad
\nu_i\in\mathcal M_{r_i}.
\]
Choose a channel $P:A\to\Delta(\{1,\dots,n\})$ implementing $\nu$ at prior $q$.  For each $i$, choose a finite
continuation channel $Q^i:A\to\Delta(Z_i)$ implementing $\nu_i$ at prior $r_i$; on actions outside the support of
$r_i$, define the rows arbitrarily.  Those rows are never reached under prior $r_i$ and do not affect $\nu_i$; moreover,
if $r_i(a)=0$, then $q(a)P(i\mid a)=m_i r_i(a)=0$, so they do not affect the compound posterior law at $q$.  The compound
channel
\[
C:=P\triangleright\{Q^i\}_{i=1}^n
\]
implements $\mu$ at prior $q$.  Let $T:\bigsqcup_{i=1}^n(\{i\}\times Z_i)\to\Delta(\{1,\dots,n\})$ be the kernel
$T(i\mid(i,z))=1$; then $CT=P$.  Clause \textup{(DP-rec)} of Condition~\textup{(P4)} applied to $C$ and $T$ gives
\[
q^0\succeq_{C\sqcup P}q^1.
\]
By Lemma~\ref{pt:lem:postsep},
\[
F_q(\mu)\ge F_q(\nu).
\]

To prove convexity, take $r,s_1,\dots,s_k\in\Delta(A)$ with
\[
r=\sum_{j=1}^k\lambda_j s_j,
\qquad
\lambda_j\ge 0,\quad \sum_j\lambda_j=1.
\]
Choose $0<\varepsilon<\min_a q(a)$ and set
\[
t:=\frac{q-\varepsilon r}{1-\varepsilon}.
\]
This belongs to $\Delta(A)$: for each $a$, $q(a)-\varepsilon r(a)\ge q(a)-\varepsilon>0$, and the coordinates sum to one.
The law
\[
\varepsilon\sum_{j=1}^k\lambda_j\delta_{s_j}+(1-\varepsilon)\delta_t
\]
is a refinement of
\[
\varepsilon\delta_r+(1-\varepsilon)\delta_t,
\]
and both belong to $\mathcal M_q$.  Refinement monotonicity and the integral representation give
\[
\varepsilon\sum_{j=1}^k\lambda_j\phi_q(s_j)+(1-\varepsilon)\phi_q(t)
\ge
\varepsilon\phi_q(r)+(1-\varepsilon)\phi_q(t).
\]
Cancelling the common term and dividing by $\varepsilon$ yields
\[
\phi_q(r)\le \sum_{j=1}^k\lambda_j\phi_q(s_j),
\]
so $\phi_q$ is convex.

Since $q$ is in the relative interior of the finite-dimensional simplex, the supporting-hyperplane theorem for finite
convex functions gives a relative subgradient $g\in\R^A$ of the continuous convex function $\phi_q$ at $q$.  Define
\[
\widetilde\phi_q(r):=\phi_q(r)-g\cdot(r-q).
\]
For every $\mu\in\mathcal M_q$,
\[
\int g\cdot(r-q)\,d\mu(r)=g\cdot(b(\mu)-q)=0,
\]
so $\widetilde\phi_q$ represents the same affine functional on $\mathcal M_q$.  The subgradient inequality gives
\[
\widetilde\phi_q(r)\ge \widetilde\phi_q(q)
\qquad\text{for all }r\in\Delta(A).
\]
Finally replace
\[
F_q(\mu)\quad\text{by}\quad F_q(\mu)-F_q(\delta_q)
\]
and replace $\widetilde\phi_q$ by
\[
\widetilde\phi_q-\widetilde\phi_q(q).
\]
Before this final constant shift, the representation with $\widetilde\phi_q$ gives
\[
F_q(\delta_q)=\int \widetilde\phi_q(r)\,d\delta_q(r)=\widetilde\phi_q(q).
\]
Thus the same constant is subtracted from both sides of the integral representation, and subtracting a constant from
$F_q$ preserves all ordinal comparisons on $\mathcal M_q$.  The final normalisation gives
\[
\phi_q(q)=F_q(\delta_q)=0,
\]
and makes $\phi_q$ non-negative on the whole simplex.  Hence
\[
F_q(\mu)=\int\phi_q(r)\,d\mu(r)\ge 0
\]
for every $\mu\in\mathcal M_q$.

Suppose now that $|A|\ge2$.  Full revelation implements $\chi_q$ and no
information implements $\delta_q$.  The strict local-orientation clause of
Condition~\textup{(P1)}, transported through the quotient comparison, gives
\[
F_q(\chi_q)>F_q(\delta_q)=0.
\]
Divide both $F_q$ and $\phi_q$ by the positive number $F_q(\chi_q)$.  This
preserves the integral representation, all comparisons, convexity, and
non-negativity, and gives $F_q(\chi_q)=1$.  These two anchor values make the
selected affine representative unique.  If $|A|=1$, every posterior law with
barycentre $q$ is $\delta_q$, so the zero representative is the only relevant
choice.
\end{proof}

For a finitely supported posterior law $\mu=\sum_i\lambda_i\delta_{r_i}$ and a bijection $\pi:A\to A'$, write
\[
\mu\pi:=\sum_i\lambda_i\delta_{r_i\pi}
\]
for the posterior law obtained by relabelling actions through $\pi$.

\begin{lemma}[Exact cross-fibre relabelling]\label{pt:lem:actionbase}
Action coarsening implies ordinal cross-fibre equivalence for bijections.  The
zero/full-revelation selection in Lemma~\ref{pt:lem:convnorm} makes that
equivalence exact: for every bijection $\pi:A\to A'$ and full-support
$q\in\Delta(A)$,
\[
F_{q\pi}(\mu\pi)=F_q(\mu)
\]
for all $\mu\in\mathcal M_q$.  In particular, no-information experiments
have the same selected value zero on every finite action set.
\end{lemma}

\begin{proof}
\emph{Ordinal equivalence.}  Let $\pi:A\to A'$ be a bijection.  View $\pi$ as the deterministic kernel
$S(a'\mid a)=\1_{a'=\pi(a)}$.  For any channel $P:A\to\Delta(X)$ and full-support $q\in\Delta(A)$, the Bayesian
pushforward satisfies $S^qP=P\pi^{-1}$ (explicitly: $(S^qP)(x\mid a')=P(x\mid\pi^{-1}(a'))$).  Clause
\textup{(DP-act)} of Condition~\textup{(P4)} gives
\[
q^0\succeq_{P\sqcup P\pi^{-1}}(q\pi)^1.
\]
For the reverse inequality, apply \textup{(DP-act)} to the bijection $\pi^{-1}:A'\to A$, the prior $q\pi\in\Delta(A')$,
and the channel $P\pi^{-1}:A'\to\Delta(X)$.  Then $(\pi^{-1})^{q\pi}(P\pi^{-1})=P$ (the inverse bijection undoes the
relabeling), and $(q\pi)\pi^{-1}=q$.  Hence
\[
(q\pi)^0\succeq_{P\pi^{-1}\sqcup P}q^1.
\]
By the irrelevant-blocks clause of Condition~\textup{(P3)}, block comparisons depend only on the pair of channels involved, not on
numeric labels.  The comparison $(q\pi)^0\succeq_{P\pi^{-1}\sqcup P}q^1$ asks whether $(q\pi,P\pi^{-1})$ is at least
as good as $(q,P)$.  Swapping the block labels (so that $P$ becomes block~0 and $P\pi^{-1}$ becomes block~1) gives
the equivalent statement $q^0\preceq_{P\sqcup P\pi^{-1}}(q\pi)^1$.  Combining with the first inequality gives
\[
q^0\sim_{P\sqcup P\pi^{-1}}(q\pi)^1.
\]
If $P$ implements $\mu\in\mathcal M_q$, then $P\pi^{-1}$ implements $\mu\pi\in\mathcal M_{q\pi}$: for every
positive-probability record $x$,
\[
m_{q\pi,P\pi^{-1}}(x)=m_{q,P}(x),
\qquad
r_x^{\,q\pi,P\pi^{-1}}=r_x^{\,q,P}\pi.
\]
Now take arbitrary $\mu,\nu\in\mathcal M_q$, and choose channels $P,Q$ implementing them at prior $q$.  Let $E$ be the
four-block environment with blocks $P$, $P\pi^{-1}$, $Q$, and $Q\pi^{-1}$.  Write $x,x',y,y'$ for the block-supported
lotteries $q$, $q\pi$, $q$, and $q\pi$ on these four blocks, respectively.  By Lemma~\ref{pt:lem:blockcoh}, the two
two-block indifferences give $x\sim_E x'$ and $y\sim_E y'$.  Hence, by transitivity of the weak order on $E$,
\[
x\succeq_E y
\quad\Longleftrightarrow\quad
x'\succeq_E y'.
\]
Applying Lemma~\ref{pt:lem:blockcoh} again to delete irrelevant blocks gives
\[
q^0\succeq_{P\sqcup Q}q^1
\quad\Longleftrightarrow\quad
(q\pi)^0\succeq_{P\pi^{-1}\sqcup Q\pi^{-1}}(q\pi)^1.
\]
By Lemma~\ref{pt:lem:postsep} on the two fibres, this is equivalent to
\[
F_q(\mu)\ge F_q(\nu)
\quad\Longleftrightarrow\quad
F_{q\pi}(\mu\pi)\ge F_{q\pi}(\nu\pi).
\]
The map $\mu\mapsto\mu\pi$ is an affine bijection from $\mathcal M_q$ to $\mathcal M_{q\pi}$, with inverse induced by
$\pi^{-1}$.  Therefore $G(\mu):=F_{q\pi}(\mu\pi)$ is a continuous affine representative of the same weak order on
$\mathcal M_q$ as $F_q$.  Herstein--Milnor uniqueness gives $G=aF_q+b$ with $a>0$.  Since
$\delta_q\pi=\delta_{q\pi}$ and Lemma~\ref{pt:lem:convnorm} gives
\[
F_q(\delta_q)=F_{q\pi}(\delta_{q\pi})=0,
\]
we have $b=0$.  If $|A|=1$, both representatives vanish identically.  If
$|A|\ge2$, relabelling sends full revelation to full revelation,
$\chi_q\pi=\chi_{q\pi}$, while the canonical selection gives
$F_q(\chi_q)=F_{q\pi}(\chi_{q\pi})=1$.  Substitution in $G=aF_q$ therefore
gives $a=1$.  Hence
\[
F_{q\pi}(\mu\pi)=F_q(\mu)
\]
for all $\mu\in\mathcal M_q$.

\emph{No-information equivalence.}  By Lemma~\ref{pt:lem:convnorm}, $F_q(\delta_q)=0$ for every full-support $q$ on every
finite action set, independently of cardinality.
\end{proof}

\begin{lemma}[Undetectable distinctions neutrality]\label{pt:lem:neutrality}
Let $S:A\to A'$ be an onto deterministic partition map, and let $P:A\to\Delta(X)$ be $S$-measurable:
\[
S(a)=S(b)\quad\Rightarrow\quad P(\cdot\mid a)=P(\cdot\mid b).
\]
Then for every full-support $q\in\Delta(A)$,
\[
q^0\sim_{P\sqcup S^qP}(qS)^1.
\]
In particular, if $G_S:A\to\Delta(A')$ reveals the block $S(a)$, then
\[
(q,G_S)\sim(qS,\Id_{A'}).
\]
\end{lemma}

\begin{proof}
Since $S$ is onto and $q$ has full support, $qS$ has full support on $A'$.  Thus the Bayesian pushforward $S^qP$ has no
zero-row ambiguity in this lemma.  View the deterministic map $S$ as its associated stochastic kernel.
One direction is exactly clause \textup{(DP-act)} of Condition~\textup{(P4)}:
\[
q^0\succeq_{P\sqcup S^qP}(qS)^1.
\]
For the reverse direction, define the Bayesian refinement kernel $R:A'\to\Delta(A)$ by
\[
R(a\mid a')=
\begin{cases}
q(a)/(qS)(a'),&S(a)=a',\\
0,&\text{otherwise}.
\end{cases}
\]
For each $a'\in A'$, $R(\cdot\mid a')$ is stochastic because
\[
\sum_{a\in A}R(a\mid a')
=
\frac{\sum_{a:S(a)=a'}q(a)}{(qS)(a')}
=1.
\]
Moreover, for every $a\in A$,
\[
((qS)R)(a)
=
\sum_{a'\in A'}(qS)(a')R(a\mid a')
=
(qS)(S(a))\frac{q(a)}{(qS)(S(a))}
=q(a)>0.
\]
Thus $(qS)R=q$.

Since $P$ is $S$-measurable, the row of $S^qP$ at each block $a'$ equals the common $P$-row on that block: if $S(a)=a'$,
then
\[
S^qP(x\mid a')
=
\frac{\sum_{b:S(b)=a'}q(b)P(x\mid b)}{(qS)(a')}
=
P(x\mid a).
\]
Now pushing $S^qP$ back through $R$ recovers $P$.  For each $a\in A$,
\[
\bigl(R^{qS}(S^qP)\bigr)(x\mid a)
=
\frac{\sum_{a'\in A'}(qS)(a')R(a\mid a')S^qP(x\mid a')}{((qS)R)(a)}
=
S^qP(x\mid S(a))
=
P(x\mid a).
\]
Therefore
\[
R^{qS}(S^qP)=P.
\]

Apply \textup{(DP-act)} to $(qS,S^qP)$ with the refinement kernel $R$:
\[
(qS)^0\succeq_{S^qP\sqcup P}q^1.
\]
By Condition~\textup{(P3)} applied to the two-block environment $P\sqcup S^qP$ with ordered pair $(1,0)$, this is equivalent
to
\[
(qS)^1\succeq_{P\sqcup S^qP}q^0.
\]
Together the two inequalities give the desired indifference.
\end{proof}

For a nonempty $B\subseteq A$ and a finite law $\mu$ whose barycentre is supported on $B$, non-negativity implies that
every atom $p$ of $\mu$ is supported on $B$.  Define its support projection by
\[
\pi_B(\mu)
:=
\sum_{p\in\operatorname{supp}(\mu)}\mu(\{p\})\,\delta_{p|_B}.
\]
It is a finite law on $\Delta(B)$ with barycentre $b(\mu)|_B$.

\begin{lemma}[Support restriction]\label{pt:lem:supprestrict}
Let $r\in\Delta(A)$ have support $B:=\operatorname{supp}(r)\subsetneq A$.  Then:
\begin{enumerate}
\item[(i)] \emph{Posterior-law reduction.}  For any channel $P:A\to\Delta(X)$, the posterior law $\mu_{r,P}$ depends only
on the rows $P(\cdot\mid b)$ for $b\in B$.  Rows $P(\cdot\mid a)$ with $a\notin B$ have zero contribution because $r(a)=0$.
\item[(ii)] \emph{Preference equivalence.}  Two channels $P,Q:A\to\Delta(X)$ that agree on $B$ (i.e.\ $P(\cdot\mid b)=Q(\cdot\mid b)$
for all $b\in B$) are indifferent at prior $r$.
\item[(iii)] \emph{Face identification.}  For every channel $P:A\to\Delta(X)$,
\[
(r,P)\sim(r|_B,P|_B)
\]
as a two-block comparison; consequently the comparison between channels $P,Q$ at prior $r$ agrees with the comparison
between $P|_B,Q|_B$ at the full-support prior $r|_B\in\Delta(B)$.  This is proved by projecting $A$ to $B$ and embedding
$B$ back into $A$, using the all-completions form of clause \textup{(DP-act)} of Condition~\textup{(P4)} for the embedding.
\item[(iv)] \emph{Representative.}  Define $F_r:\mathcal M_r\to\R$ by $F_r(\mu):=F_{r|_B}^B(\pi_B(\mu))$, where $\pi_B$
projects posterior laws from $\Delta(A)$ to $\Delta(B)$ and $F_{r|_B}^B$ is the Lemma~\ref{pt:lem:postsep} representative
on $\Delta(B)$.  This $F_r$ represents continuation comparisons at prior $r$ in action set $A$.
\end{enumerate}
\end{lemma}

\begin{proof}
In this proof, $(s,R)\succeq(t,T)$ abbreviates the block comparison
$s^0\succeq_{R\sqcup T}t^1$, and $\sim$ denotes both weak inequalities.  These comparisons involve different priors and
different action sets; they chain by transitivity in Lemma~\ref{pt:lem:blockcoh}(a), which is stated across such pairs.

Part~(i): The marginal probability $m_{r,P}(x)=\sum_a r(a)P(x\mid a)$ and the Bayes numerator $r(a)P(x\mid a)$ are both
zero for $a\notin B$.  Hence $\mu_{r,P}$ is determined by $\{P(\cdot\mid b):b\in B\}$.

Part~(ii): Let $P$ and $\widetilde P$ be channels on $A$ that agree on $B$.  We first show
$(r,P)\sim(r,\widetilde P)$.  Fix $b_0\in B$, and let $S:A\to B$ be the deterministic map with $S(b)=b$ for $b\in B$
and $S(a)=b_0$ for $a\notin B$.  Then $rS=r|_B$ and the Bayesian pushforward $S^rP$ is $P|_B$ on $B$.  Clause \textup{(DP-act)}
gives
\[
(r,P)\succeq(r|_B,P|_B).
\]
Let $I:B\hookrightarrow A$ be the inclusion.  Then $(r|_B)I=r$, and the Bayesian pushforward $I^{r|_B}(P|_B)$ is
pinned down on $B$ and unconstrained on $A\setminus B$.  By the all-completions convention in \textup{(DP-act)}, choose
the completion to be $\widetilde P$.  Clause \textup{(DP-act)} gives
\[
(r|_B,P|_B)\succeq(r,\widetilde P).
\]
Thus $(r,P)\succeq(r,\widetilde P)$.  Applying the same projection/inclusion argument with $\widetilde P$ as the original
channel and choosing $P$ as the inclusion completion gives $(r,\widetilde P)\succeq(r,P)$.  Hence
$(r,P)\sim(r,\widetilde P)$.

If $P,Q$ agree on $B$, choose a common extension $\widetilde P=\widetilde Q$ of this common restriction to all of $A$
(for instance, use one fixed off-support row).  The preceding paragraph gives
$(r,P)\sim(r,\widetilde P)$ and $(r,Q)\sim(r,\widetilde Q)$; since $\widetilde P=\widetilde Q$, common-environment
transitivity gives $(r,P)\sim(r,Q)$.

Part~(iii): We first prove $(r,P)\sim(r|_B,P|_B)$ for every channel $P:A\to\Delta(X)$.

\emph{Forward direction (\textup{(DP-act)} projection).}  Let $S:A\to B$ map all $a\notin B$ to some fixed $b_0\in B$.  The coarsened
prior is $rS=r|_B$ (since $r(a)=0$ for $a\notin B$), and the Bayesian pushforward $S^rP$ equals $P|_B$ on $B$.  By
\textup{(DP-act)},
\[
(r,P)\succeq(r|_B,P|_B).
\]

\emph{Reverse direction (\textup{(DP-act)} inclusion with completion).}  Let $I:B\hookrightarrow A$ be the inclusion.  Then
$(r|_B)I=r$.  The Bayesian pushforward $I^{r|_B}(P|_B)$ is pinned down on $B$ and arbitrary on $A\setminus B$.  Choose
the \textup{(DP-act)} completion to be any extension $\widetilde P:A\to\Delta(X)$ of $P|_B$.  Then
\[
(r|_B,P|_B)\succeq(r,\widetilde P).
\]
By Part~(ii), $(r,\widetilde P)\sim(r,P)$ because $\widetilde P$ and $P$ agree on $B=\operatorname{supp}(r)$.  Hence
$(r|_B,P|_B)\succeq(r,P)$.

Together the two directions give $(r,P)\sim(r|_B,P|_B)$.

\emph{Comparison equivalence.}  For any $P,Q$: $(r,P)\succeq(r,Q)$ iff (by the above indifferences $(r,P)\sim(r|_B,P|_B)$,
$(r,Q)\sim(r|_B,Q|_B)$, and Condition~\textup{(P1)} transitivity) $(r|_B,P|_B)\succeq(r|_B,Q|_B)$.

Part~(iv): If $\mu\in\mathcal M_r$, every posterior in the finite support of $\mu$ is supported on $B$; otherwise the
barycentre $r$ would assign positive mass to $A\setminus B$.  Hence $\pi_B(\mu)\in\mathcal M_{r|_B}$ is well defined.
By (iii), continuation comparisons at $r$ are isomorphic to those at $r|_B$.  The Lemma~\ref{pt:lem:postsep} representative
$F_{r|_B}^B$ on $\Delta(B)$ therefore induces a representative on $\mathcal M_r$ via $\pi_B$.
\end{proof}

\paragraph{Stage 3: product transport and the cross-prior bridge.}
\addcontentsline{toc}{subsubsection}{Stage 3: product transport and the cross-prior bridge}
An independent background is first shown to be ordinally inert.  Separate
affinity then yields the most general compatible product expression,
$x+y+\kappa xy$, after a coherent positive gauge.  Crucially,
Lemma~\ref{pt:lem:blockbridge} is proved before any coarse-revelation value is
used: product structure and face scales alone do not identify cross-prior
block comparisons.

\begin{lemma}[Derived background inertness]\label{pt:lem:background}
Assume Conditions~\textup{(P1)} and \textup{(P3)}--\textup{(P5)}.  Let $A_1,A_2$ be finite action sets and let
$q_1\in\Delta(A_1)$ and $q_2\in\Delta(A_2)$ have full support.  Then, for all channels $P_1,Q_1$ on $A_1$ and every
background channel $R_2$ on $A_2$,
\begin{equation}
(q_1\otimes q_2)^0
\succeq_{(P_1\otimes R_2)\sqcup(Q_1\otimes R_2)}
(q_1\otimes q_2)^1
\quad\Longleftrightarrow\quad
q_1^0\succeq_{P_1\sqcup Q_1}q_1^1,
\tag{BI}\label{pt:eq:backgroundinert}
\end{equation}
and symmetrically for comparisons in the second component with a common first-component background.  In particular the
left-hand comparison does not depend on the common background: for all backgrounds $R_2,S_2$ on $A_2$,
\[
(q_1\otimes q_2)^0
\succeq_{(P_1\otimes R_2)\sqcup(Q_1\otimes R_2)}
(q_1\otimes q_2)^1
\quad\Longleftrightarrow\quad
(q_1\otimes q_2)^0
\succeq_{(P_1\otimes S_2)\sqcup(Q_1\otimes S_2)}
(q_1\otimes q_2)^1.
\]
\end{lemma}

\begin{proof}
Throughout, for a prior $s$ on an action set $A$ and channels $V,W$ on $A$ write $V\succeq_sW$ for the block comparison
$s^0\succeq_{V\sqcup W}s^1$, and $(s,V)\sim(t,W)$ for the two-block indifference between the pairs.  By
Lemma~\ref{pt:lem:blockcoh}(a) these comparisons form a weak order for each fixed prior and chain transitively across
pairs with different priors and action sets.  By Lemma~\ref{pt:lem:plsuff}, at a full-support prior a channel may be replaced by any channel with the
same posterior law, in weak and in strict comparisons alike; we call this \emph{substitution}.

\emph{Preliminary reduction: a common foreground alphabet.}
Condition~\textup{(P5)} requires the two continuation channels used on a given branch to share a record alphabet.  If
$P_1:A_1\to\Delta(X)$ and $Q_1:A_1\to\Delta(X')$, replace them by their zero-padded versions on $X\sqcup X'$, as in
Step~2 of Lemma~\ref{pt:lem:publiccoin}.  Padding adds only null records, so it changes neither $\mu_{q_1,\cdot}$ nor
$\mu_{q_1\otimes q_2,\cdot\otimes R_2}$; both sides of~\eqref{pt:eq:backgroundinert} are therefore unaffected by
substitution at the full-support priors $q_1$ and $q_1\otimes q_2$.  Assume from now on that $P_1$ and $Q_1$ share the
record alphabet $X_1$.

For a channel $V:A_1\to\Delta(X_1)$ and a finite set $A_2$, write $\widetilde V$ for its lift to the product action set,
\[
\widetilde V(z\mid(a_1,a_2)):=V(z\mid a_1),
\]
and let $\widetilde R_2(x_2\mid(a_1,a_2)):=R_2(x_2\mid a_2)$ be the lift of the background.

\emph{Step 1: a product channel is a two-stage channel with background first stage.}
Take $\widetilde R_2$ as first stage on the action set $A_1\times A_2$ and the constant continuation profile
$\widetilde V$ in every branch, so that every branch uses the same continuation alphabet $X_1$ and the compound record
set $\bigsqcup_{x_2\in X_2}(\{x_2\}\times X_1)$ is $X_2\times X_1$.  By the definition of two-stage composition,
\[
\bigl(\widetilde R_2\triangleright\{\widetilde V\}_{x_2}\bigr)(x_2,z\mid(a_1,a_2))
=R_2(x_2\mid a_2)\,V(z\mid a_1)
=(V\otimes R_2)(z,x_2\mid (a_1,a_2)).
\]
The two channels therefore differ only by the bijection $(x_2,z)\mapsto(z,x_2)$ of record labels, which leaves every
record probability and every posterior unchanged.  Hence they induce the same posterior law at $q_1\otimes q_2$, and
substitution makes them interchangeable in every comparison at that prior.  Applying this to $V=P_1$ and $V=Q_1$, the
left-hand side of~\eqref{pt:eq:backgroundinert} is equivalent to
\begin{equation}
\widetilde R_2\triangleright\{\widetilde P_1\}
\;\succeq_{q_1\otimes q_2}\;
\widetilde R_2\triangleright\{\widetilde Q_1\},
\tag{$\ast$}\label{pt:eq:bgcompound}
\end{equation}
and likewise for the strict comparisons.

\emph{Step 2: every reached branch posterior is a product with first marginal $q_1$.}
Write $q:=q_1\otimes q_2$, $m(x_2):=m_{q,\widetilde R_2}(x_2)$ and, for $m(x_2)>0$, $r_{x_2}:=r_{x_2}^{\,q,\widetilde R_2}$.
Since the rows of $\widetilde R_2$ do not depend on $a_1$,
\[
m(x_2)=\sum_{a_2\in A_2}q_2(a_2)R_2(x_2\mid a_2)=m_{q_2,R_2}(x_2),
\]
and Bayes' rule gives
\[
r_{x_2}(a_1,a_2)
=\frac{q_1(a_1)q_2(a_2)R_2(x_2\mid a_2)}{m(x_2)}
=q_1(a_1)\,q_2^{x_2}(a_2),
\qquad
q_2^{x_2}:=r_{x_2}^{\,q_2,R_2}.
\]
Thus $r_{x_2}=q_1\otimes q_2^{x_2}$: the background signal may change beliefs about $a_2$, but in every reached branch
the two components stay independent and the first marginal is still $q_1$.  At least one branch has positive
probability.

\emph{Step 3: a lifted foreground channel is evaluated as on its own component.}
We claim that for every $s_2\in\Delta(A_2)$ and every channel $V:A_1\to\Delta(X_1)$,
\begin{equation}
(q_1\otimes s_2,\widetilde V)\sim(q_1,V).
\tag{$\dagger$}\label{pt:eq:liftneutral}
\end{equation}
First suppose $s_2$ has full support, so that $q_1\otimes s_2$ has full support on $A_1\times A_2$.  Let
$S:A_1\times A_2\to A_1$ be the coordinate projection, an onto deterministic map, and note that $\widetilde V$ is
$S$-measurable, since $S(a)=S(b)$ implies $\widetilde V(\cdot\mid a)=\widetilde V(\cdot\mid b)$.  Moreover
$(q_1\otimes s_2)S=q_1$ and, for every $a_1$,
\[
\bigl(S^{q_1\otimes s_2}\widetilde V\bigr)(z\mid a_1)
=\frac{\sum_{a_2\in A_2}q_1(a_1)s_2(a_2)V(z\mid a_1)}{q_1(a_1)}
=V(z\mid a_1),
\]
so the Bayesian pushforward of $\widetilde V$ under $S$ is $V$ itself.  Undetectable-distinctions neutrality
(Lemma~\ref{pt:lem:neutrality}) gives~\eqref{pt:eq:liftneutral}.

If $s_2$ is not full support, let $B_2:=\operatorname{supp}(s_2)$, so that $B:=A_1\times B_2$ is the support of
$q_1\otimes s_2$ and $B\subsetneq A_1\times A_2$.  The set $B$ \emph{is} the product action set $A_1\times B_2$, so no
relabelling is involved: the restricted prior is
$q_1\otimes(s_2|_{B_2})$ and the restricted channel $\widetilde V|_B$ is the lift of $V$ to $A_1\times B_2$.  Face
identification (Lemma~\ref{pt:lem:supprestrict}(iii)) gives
$(q_1\otimes s_2,\widetilde V)\sim(q_1\otimes s_2|_{B_2},\widetilde V|_B)$, and the full-support case just proved,
applied on the action set $A_1\times B_2$, gives $(q_1\otimes s_2|_{B_2},\widetilde V|_B)\sim(q_1,V)$.  Transitivity in a
common block environment yields~\eqref{pt:eq:liftneutral} in general.

Fix a branch with $m(x_2)>0$ and apply~\eqref{pt:eq:liftneutral} at $s_2=q_2^{x_2}$ to $V=P_1$ and to $V=Q_1$, using
$r_{x_2}=q_1\otimes q_2^{x_2}$ from Step~2.  Place the four pairs
$(r_{x_2},\widetilde P_1)$, $(q_1,P_1)$, $(r_{x_2},\widetilde Q_1)$, $(q_1,Q_1)$ as blocks of one finite block environment
(Condition~\textup{(P3)}, Lemma~\ref{pt:lem:blockcoh}).  The two indifferences and transitivity of the weak order in that
environment (Condition~\textup{(P1)}) give
\begin{equation}
\widetilde P_1\succeq_{r_{x_2}}\widetilde Q_1
\quad\Longleftrightarrow\quad
P_1\succeq_{q_1}Q_1,
\tag{$\ddagger$}\label{pt:eq:branchreduce}
\end{equation}
and, since the relation is complete, the same equivalence for the strict comparisons.  In particular the branch
comparison is the same in every reached branch and does not involve $R_2$.

\emph{Step 4: forward implication.}
Assume $P_1\succeq_{q_1}Q_1$.  By~\eqref{pt:eq:branchreduce}, the branch hypothesis
$r_{x_2}^0\succeq_{\widetilde P_1\sqcup\widetilde Q_1}r_{x_2}^1$ of Condition~\textup{(P5)} holds for every branch with
$m(x_2)>0$, with the common first stage $\widetilde R_2$ and the two constant continuation profiles.
Condition~\textup{(P5)} yields~\eqref{pt:eq:bgcompound}, hence the left-hand side of~\eqref{pt:eq:backgroundinert} by Step~1.

\emph{Step 5: converse, via the strictness clause.}
Assume the left-hand side of~\eqref{pt:eq:backgroundinert}, so that~\eqref{pt:eq:bgcompound} holds by Step~1, and suppose
$P_1\succeq_{q_1}Q_1$ fails.  By completeness and the reverse-block equivalence of Lemma~\ref{pt:lem:blockcoh},
$Q_1\succ_{q_1}P_1$, so by the strict form of~\eqref{pt:eq:branchreduce} every reached branch satisfies
$\widetilde Q_1\succ_{r_{x_2}}\widetilde P_1$.  At least one such branch exists and has positive probability, so the
strictness clause of Condition~\textup{(P5)}, applied with the two profiles exchanged, gives
\[
\widetilde R_2\triangleright\{\widetilde Q_1\}
\;\succ_{q_1\otimes q_2}\;
\widetilde R_2\triangleright\{\widetilde P_1\},
\]
contradicting~\eqref{pt:eq:bgcompound} because $\succeq_{q_1\otimes q_2}$ is a weak order
(Lemma~\ref{pt:lem:blockcoh}).  Hence $P_1\succeq_{q_1}Q_1$, proving~\eqref{pt:eq:backgroundinert}.

The second-component clause is the same argument with the roles of the two coordinates exchanged: the common
first-component background is used as the first stage and the second-component channels as the continuations.  Finally,
since both product comparisons in the last display of the lemma are equivalent to the single background-free comparison
$P_1\succeq_{q_1}Q_1$, they are equivalent to each other.
\end{proof}

The lemma shows that separability across independent components is not an independent behavioural commitment.  A product
channel is a two-stage channel whose first stage is the background experiment: it can teach the observer about the
background action, but it leaves the foreground component independent of it and its marginal unchanged, so in every
reached branch the continuation comparison is the foreground comparison itself.  Branchwise monotonicity then transmits
that single comparison to the compound, and its strictness clause transmits it back.  The derived statement is stronger
than mere invariance to changing the background: it identifies the product comparison with the comparison of the
foreground channels alone, which is the form used in Lemma~\ref{pt:lem:coherentnorm} below.

For finite posterior laws $\mu=\sum_i m_i\delta_{r_i}\in\mathsf X_A$ and
$\nu=\sum_j n_j\delta_{s_j}\in\mathsf X_B$, define
\[
\mu\otimes\nu
:=
\sum_{i,j}m_i n_j\,\delta_{r_i\otimes s_j}.
\]
Bayes' rule then gives the product factorisation
\begin{equation}
\mu_{q_1\otimes q_2,\,P_1\otimes P_2}
=
\mu_{q_1,P_1}\otimes\mu_{q_2,P_2}.
\tag{PL}\label{pt:eq:productlaw}
\end{equation}

\begin{lemma}[Coherent product quasi-additivity]\label{pt:lem:coherentnorm}
There is a choice of positive rescalings of the zero-normalised fibrewise representatives $F_q$, still denoted $F_q$,
and a scalar $\kappa\in\R$ such that, for all full-support priors $p\in\Delta(A)$ and $r\in\Delta(B)$ and all
$\mu\in\mathcal M_p$, $\nu\in\mathcal M_r$,
\begin{equation}
F_{p\otimes r}(\mu\otimes\nu)
=F_p(\mu)+F_r(\nu)+\kappa F_p(\mu)F_r(\nu).
\tag{QA}\label{pt:eq:QA}
\end{equation}
If one factor is a singleton, the same formula holds with its identically zero representative.  Moreover, for every
nondegenerate $r$ and every $\nu\in\mathcal M_r$,
\begin{equation}
1+\kappa F_r(\nu)>0.
\tag{POS}\label{pt:eq:QAslope}
\end{equation}
\end{lemma}

\begin{proof}
All representatives in this proof are zero-normalised: $F_q(\delta_q)=0$.  The proof first treats nondegenerate action
sets, meaning action sets with at least two elements.

\emph{Step 1: product-coordinate independence.}
Fix full-support priors $p$ and $r$ on nondegenerate finite action sets and set
\[
L_{p,r}(\mu,\nu):=F_{p\otimes r}(\mu\otimes\nu).
\]
This map is continuous and separately affine.  To identify a first-coordinate slice, let $P,Q$ implement
$\mu,\mu'\in\mathcal M_p$ and let $R$ implement $\nu\in\mathcal M_r$.  By the product factorisation
\eqref{pt:eq:productlaw}, $P\otimes R$ implements $\mu\otimes\nu$ and $Q\otimes R$ implements $\mu'\otimes\nu$ at the
full-support prior $p\otimes r$.  Hence Lemma~\ref{pt:lem:postsep}, applied on the fibre $\mathcal M_{p\otimes r}$ and on
the fibre $\mathcal M_p$, turns background inertness
\eqref{pt:eq:backgroundinert} of Lemma~\ref{pt:lem:background} into
\[
L_{p,r}(\mu,\nu)\ge L_{p,r}(\mu',\nu)
\quad\Longleftrightarrow\quad
F_p(\mu)\ge F_p(\mu').
\]
The symmetric clause of Lemma~\ref{pt:lem:background} identifies every second-coordinate slice with the $F_r$-order.  Local non-triviality, the
fibrewise representation, and the zero normalisation from Lemma~\ref{pt:lem:convnorm} make both slice orders nonconstant:
$F_p(\chi_p)>0$ and $F_r(\chi_r)>0$.

\emph{Step 2: the pair-specific bilinear form.}
Put $x(\mu):=F_p(\mu)$ and $y(\nu):=F_r(\nu)$.  For fixed $\nu$, affine-utility uniqueness gives
\[
L_{p,r}(\mu,\nu)=\alpha(\nu)x(\mu)+\gamma(\nu),
\qquad \alpha(\nu)>0.
\]
Taking $\mu=\delta_p$ shows $\gamma(\nu)=L_{p,r}(\delta_p,\nu)$.  This is a continuous affine representative of the
$F_r$-order and vanishes at $\delta_r$, so
\[
\gamma(\nu)=B_{p,r}y(\nu)
\]
for some $B_{p,r}>0$.  Likewise, the slice at $\nu=\delta_r$ gives
\[
L_{p,r}(\mu,\delta_r)=A_{p,r}x(\mu)
\]
for some $A_{p,r}>0$.

Let $\bar\mu:=\chi_p$ and $h_p:=F_p(\chi_p)>0$.  The function
$\nu\mapsto L_{p,r}(\bar\mu,\nu)$ is again a continuous affine representative of the $F_r$-order.  Hence there are
$D_{p,r}>0$ and a constant $E_{p,r}$ such that
\[
L_{p,r}(\bar\mu,\nu)=D_{p,r}y(\nu)+E_{p,r}.
\]
At $\nu=\delta_r$, $E_{p,r}=A_{p,r}h_p$.  Comparing this identity with
$L_{p,r}(\bar\mu,\nu)=\alpha(\nu)h_p+B_{p,r}y(\nu)$ gives
\[
\alpha(\nu)=A_{p,r}+C_{p,r}y(\nu),
\qquad
C_{p,r}:=\frac{D_{p,r}-B_{p,r}}{h_p}.
\]
Therefore
\begin{equation}
L_{p,r}(\mu,\nu)
=A_{p,r}x(\mu)+B_{p,r}y(\nu)+C_{p,r}x(\mu)y(\nu).
\tag{BIL}\label{pt:eq:pairbilinear}
\end{equation}
The coordinate-slice slopes are positive:
\[
A_{p,r}+C_{p,r}y(\nu)>0,
\qquad
B_{p,r}+C_{p,r}x(\mu)>0.
\]

\emph{Step 3: normalising the two linear coefficients.}
Let
\[
\alpha:((A\times B)\times C)\longrightarrow A\times(B\times C),
\qquad \alpha((a,b),c)=(a,(b,c)),
\]
be the canonical associating bijection.  It carries the left-grouped product
prior and posterior law to the right-grouped ones.  Exact cross-fibre
relabelling (Lemma~\ref{pt:lem:actionbase}) therefore makes the two evaluations
equal, with no unknown affine factor.  Expanding this common value under the
two groupings and comparing the coefficients of $x$, $y$, and $z$ gives
\begin{align}
A_{p\otimes r,s}A_{p,r}&=A_{p,r\otimes s}, \tag{C1}\label{pt:eq:Acoc1}\\
A_{p\otimes r,s}B_{p,r}&=B_{p,r\otimes s}A_{r,s}, \tag{C2}\label{pt:eq:Acoc2}\\
B_{p\otimes r,s}&=B_{p,r\otimes s}B_{r,s}. \tag{C3}\label{pt:eq:Acoc3}
\end{align}
Coefficient comparison is legitimate because the mixtures $(1-t)\delta_p+t\chi_p$ make $x$ range over a nondegenerate
interval, and similarly in the other coordinates.

The coordinate swap is another bijection, so exact cross-fibre relabelling gives
\[
F_{r\otimes p}(\nu\otimes\mu)=F_{p\otimes r}(\mu\otimes\nu).
\]
Comparing coefficients in~\eqref{pt:eq:pairbilinear},
\[
B_{r,p}=A_{p,r},
\qquad
A_{r,p}=B_{p,r},
\qquad
C_{r,p}=C_{p,r}.
\]
Thus, with $\rho(p,r):=B_{p,r}/A_{p,r}$,
\begin{equation}
\rho(r,p)=\rho(p,r)^{-1}.
\tag{SW}\label{pt:eq:rhoswap}
\end{equation}

Divide~\eqref{pt:eq:Acoc2} by~\eqref{pt:eq:Acoc1}:
\[
\rho(p,r)=\rho(p,r\otimes s)A_{r,s}.
\]
Fix a nondegenerate reference prior $q_0$ on a finite action set $A_0$ and put $g(r):=\rho(q_0,r)$.  Then
\begin{equation}
A_{r,s}=\frac{g(r)}{g(r\otimes s)}.
\tag{G}\label{pt:eq:gaugeA}
\end{equation}
The coefficients $A_{p,r},B_{p,r},C_{p,r}$ are unchanged when either prior is
transported by a bijection.  Indeed, transport the product law by the product
bijection, use Lemma~\ref{pt:lem:actionbase} on the two input fibres and the
product fibre, and compare the unique coefficients in
\eqref{pt:eq:pairbilinear}.  Hence $\rho$, and therefore $g$, is invariant
under relabelling.  In particular it takes the same value on the two
associations of a product prior.

Rescale every representative by $F_p^*:=g(p)F_p$.  Under this positive rescaling,
\[
A^*_{p,r}=A_{p,r}\frac{g(p\otimes r)}{g(p)},
\quad
B^*_{p,r}=B_{p,r}\frac{g(p\otimes r)}{g(r)},
\quad
C^*_{p,r}=C_{p,r}\frac{g(p\otimes r)}{g(p)g(r)}.
\]
Because $g$ is relabelling-invariant, these rescaled representatives retain
the exact relabelling equality of Lemma~\ref{pt:lem:actionbase}.  Thus the
associator and swap comparisons used below still carry no scalar defect.
Equation~\eqref{pt:eq:gaugeA} gives $A^*_{p,r}=1$ for every pair.  The
transformed swap identity gives $B^*_{p,r}=A^*_{r,p}=1$.  Thus, after
the positive gauge rescaling, $A_{p,r}=B_{p,r}=1$ for every nondegenerate pair
(we now drop the stars).  Positivity of the slice
slopes is preserved because the rescaling factors are positive.

For a boundary prior $r\in\Delta(A)$ with $B=\operatorname{supp}(r)$, redefine
the post-gauge boundary representative by exact support transport,
\[
F_r^A(\mu):=F_{r|_B}^B\bigl(\pi_B(\mu)\bigr),
\]
where the representative on the right is the newly gauged full-support one;
singleton support faces remain identically zero.  Thus the support identity of
Lemma~\ref{pt:lem:supprestrict} remains definitional after the gauge change.

\emph{Step 4: the interaction coefficient is universal.}
Write the remaining coefficient as $\kappa_{p,r}$.  Comparing the $xy$, $xz$, $yz$, and $xyz$ coefficients in the two
associative expansions gives
\begin{align}
\kappa_{p,r}&=\kappa_{p,r\otimes s}, \tag{K1}\label{pt:eq:kappa1}\\
\kappa_{p\otimes r,s}&=\kappa_{p,r\otimes s}, \tag{K2}\label{pt:eq:kappa2}\\
\kappa_{p\otimes r,s}&=\kappa_{r,s}, \tag{K3}\label{pt:eq:kappa3}\\
\kappa_{p\otimes r,s}\kappa_{p,r}&=\kappa_{p,r\otimes s}\kappa_{r,s}. \tag{K4}\label{pt:eq:kappa4}
\end{align}
The last equation follows from the first three, but it records the full coefficient comparison.  After $A=B=1$, exact
invariance under the coordinate swap and comparison of the bilinear coefficient give
\begin{equation}
\kappa_{p,r}=\kappa_{r,p}.
\tag{KS}\label{pt:eq:kappasym}
\end{equation}
Using~\eqref{pt:eq:kappa1}--\eqref{pt:eq:kappa3} with the reference prior $q_0$,
\[
\kappa_{p,r}
=\kappa_{p,r\otimes q_0}
=\kappa_{p\otimes r,q_0}
=\kappa_{r,q_0}.
\]
Hence $\kappa_{p,r}$ is independent of its first argument.  By symmetry it is also independent of its second argument.
Denote the common value by $\kappa$.  This proves~\eqref{pt:eq:QA} for nondegenerate factors.

The positive slice-slope condition now reads
\[
1+\kappa F_r(\nu)>0,
\]
which is~\eqref{pt:eq:QAslope}.

\emph{Step 5: singleton factors.}
Let $q_*=\delta_*$ be the prior on a singleton.  Its zero-normalised representative is identically zero.  If $p$ is
nondegenerate, the canonical bijection $A\to A\times\{*\}$ and exact
relabeling give
\[
F_{p\otimes q_*}(\mu\otimes\delta_{q_*})=F_p(\mu).
\]
The product gauge preserves this equality because $p\otimes q_*$ is a
relabelling of $p$.  Since $F_{q_*}=0$, equation~\eqref{pt:eq:QA} follows
immediately.  The reversed order is identical, and the two-singleton case is
$0=0$.
\end{proof}

\begin{corollary}[Exact relabelling after the product gauge]\label{pt:cor:permutationinvariance}
For every bijection $\pi:A\to A'$ and every full-support $q\in\Delta(A)$,
\[
F_{q\pi}(\mu\pi)=F_q(\mu)
\qquad(\mu\in\mathcal M_q).
\]
\end{corollary}

\begin{proof}
Before the product gauge this is Lemma~\ref{pt:lem:actionbase}.  The proof of
Lemma~\ref{pt:lem:coherentnorm} shows that the gauge factor $g(q)$ is
unchanged by every bijection, so multiplication by that factor preserves the
equality.  The singleton case is immediate because both representatives
vanish.
\end{proof}

\begin{lemma}[Cross-prior block representation]\label{pt:lem:blockbridge}
For arbitrary finite priors $q\in\Delta(A)$ and $r\in\Delta(B)$ and channels
$P:A\to\Delta(X)$ and $Q:B\to\Delta(Y)$,
\begin{equation}
q^0\succeq_{P\sqcup Q}r^1
\quad\Longleftrightarrow\quad
F_q(\mu_{q,P})\ge F_r(\mu_{r,Q}).
\tag{BC}\label{pt:eq:blockcomp}
\end{equation}
Boundary representatives are interpreted by Lemma~\ref{pt:lem:supprestrict}.
\end{lemma}

\begin{proof}
First suppose $q$ and $r$ have full support, allowing singleton action sets.  At prior $q\otimes r$, quasi-additivity and
zero normalisation give
\[
F_{q\otimes r}(\mu_{q,P}\otimes\delta_r)=F_q(\mu_{q,P}),
\qquad
F_{q\otimes r}(\delta_q\otimes\mu_{r,Q})=F_r(\mu_{r,Q}).
\]
These values are in the same fibre $\mathcal M_{q\otimes r}$, so Lemma~\ref{pt:lem:postsep} gives
\[
(q\otimes r)^0\succeq_{(P\otimes U_B)\sqcup(U_A\otimes Q)}(q\otimes r)^1
\quad\Longleftrightarrow\quad
F_q(\mu_{q,P})\ge F_r(\mu_{r,Q}).
\]

It remains to transfer the product-lifted comparison back to the original two blocks.  For the $A$-side pair,
Clause \textup{(DP-act)} of Condition~\textup{(P4)} applied to the projection $\pi_A:A\times B\to A$ gives
\[
(q\otimes r,P\otimes U_B)\succeq(q,P)
\]
in block comparisons.  Conversely, let $T:A\to\Delta(A\times B)$ be
\[
T((a,b)\mid a')=\1_{\{a=a'\}}r(b).
\]
Then $qT=q\otimes r$, and the Bayesian pushforward $T^qP$ has row $P(\cdot\mid a)$ at every $(a,b)$.  Thus
$T^qP$ is $P\otimes U_B$ after the canonical one-record identification of the $B$-experiment, and \textup{(DP-act)}
gives the reverse comparison.  Equivalently, if the harmless record identification is not made, Lemma~\ref{pt:lem:plsuff}
identifies $T^qP$ and $P\otimes U_B$ because they induce the same posterior law at $q\otimes r$.  Hence
$(q\otimes r,P\otimes U_B)$ and $(q,P)$ are indifferent in every relevant block comparison.  The same
projection/embedding argument identifies $(q\otimes r,U_A\otimes Q)$ with $(r,Q)$.

Place the four pairs $(q,P)$, $(q\otimes r,P\otimes U_B)$, $(r,Q)$, and $(q\otimes r,U_A\otimes Q)$ as blocks of one
finite block environment.  Lemma~\ref{pt:lem:blockcoh} and transitivity transfer the two pairwise indifferences to the
two-block comparison, proving~\eqref{pt:eq:blockcomp} for full-support priors.

For arbitrary priors, let $A_+:=\operatorname{supp}(q)$ and $B_+:=\operatorname{supp}(r)$.  Lemma~\ref{pt:lem:supprestrict}
identifies $(q,P)$ with $(q|_{A_+},P|_{A_+})$, identifies $(r,Q)$ with $(r|_{B_+},Q|_{B_+})$, and identifies the two
values $F_q(\mu_{q,P})$ and $F_r(\mu_{r,Q})$ with their support-face values.  Applying the full-support bridge on
$A_+$ and $B_+$ and pulling back by Lemma~\ref{pt:lem:supprestrict} proves the claim.  Singleton supports use the same
zero-representative convention as Lemma~\ref{pt:lem:coherentnorm}.
\end{proof}

\paragraph{Stage 4: branch coefficients, cocycle, and support faces.}
\addcontentsline{toc}{subsubsection}{Stage 4: branch coefficients, cocycle, and support faces}
Insertion into a reached branch is affine.  Its linear part is evaluated on
finite atomic signed posterior laws; the spanning argument shows that the
resulting positive coefficient depends only on the prior and reached
posterior.  Path independence gives the cocycle and the normalized chain
rule.  The subsequent cardinal correction aligns auxiliary chain scales on
nested faces; it does not add a behavioural assumption or change the ordinal
representative.

\begin{lemma}[Branch aggregation]\label{pt:lem:branchagg}
Let $P_1:A\to\Delta(X_1)$ be a first-stage channel at full-support $q$, with branch probabilities $m(x)$ and posteriors
$r_x$.  There are positive branch coefficients $\beta(q,r_x)$, depending only on $q$ and the reached posterior $r_x$,
such that for every continuation profile $\{Q^x\}_{x:m(x)>0}$,
\[
F_q(\mu_{q,P_1\triangleright\{Q^x\}})
=
F_q(\mu_{q,P_1})+
\sum_{x:m(x)>0}m(x)\,\beta(q,r_x)\,
F_{r_x}(\mu_{r_x,Q^x}).
\]
If $r_x$ has full support, $\beta(q,r_x)$ is the full-support tangent coefficient.  If
$1<|\operatorname{supp}(r_x)|<|A|$, it is the full-to-face coefficient.  If $|\operatorname{supp}(r_x)|=1$, the value
$F_{r_x}$ is identically zero and $\beta(q,r_x)$ is an arbitrary positive convention.
\end{lemma}

\begin{proof}
Fix the first-stage channel $P_1$ and its branch structure $(m(x),r_x)_{x:m(x)>0}$.  For any profile of continuation
channels $\{Q^x\}$, the compound posterior law at prior $q$ is
\[
\mu_{q,P_1\triangleright\{Q^x\}}
=
\sum_{x:m(x)>0}m(x)\,\mu_{r_x,Q^x}.
\]

\emph{Setup: linear part of affine functional.}  Since $F_q$ is affine on $\mathcal M_q$ (Lemma~\ref{pt:lem:postsep}),
it extends uniquely to an affine functional on the affine hull of $\mathcal M_q$.  Write $L_q$ for the linear part:
for any $\mu,\mu'\in\mathcal M_q$,
\[
F_q(\mu)-F_q(\mu')=L_q(\mu-\mu').
\]

\emph{Branch-affine structure.}  Fix a positive-probability branch $\bar x$ and fix continuation channels
$H^x:A\to\Delta(Z^x)$ in all other branches $x\ne\bar x$.  Let $r:=r_{\bar x}$ and $m:=m(\bar x)$.  If
$\operatorname{supp}(r)$ is a singleton, then $\mathcal M_r=\{\delta_r\}$ and $F_r(\delta_r)=0$ by the singleton
zero-normalisation convention; such branches contribute zero and their coefficient convention is recorded in
Case~1 below.  For the branch-affine argument in this paragraph, assume $|\operatorname{supp}(r)|\ge2$.  If $r$ is a
boundary posterior, $F_r$ and continuation comparisons at $r$ are understood through Lemma~\ref{pt:lem:supprestrict}.

Let
\[
\rho_H:=\sum_{x\ne\bar x}m(x)\,\mu_{r_x,H^x}.
\]
For any $\nu\in\mathcal M_r$, the aggregate posterior law produced by using branch law $\nu$ in branch $\bar x$ and
the fixed continuations elsewhere is
\[
T_H(\nu):=m\nu+\rho_H\in\mathcal M_q.
\]
Define $g:\mathcal M_r\to\R$ by
\[
g(\nu):=F_q(T_H(\nu)).
\]
Thus $g$ depends only on $\nu$, and $g=F_q\circ T_H$ is continuous and affine.

We now show that $g$ and $F_r$ represent the same weak order on $\mathcal M_r$.  For $\nu,\nu'\in\mathcal M_r$, choose
finite continuation channels implementing them at prior $r$; when $r$ is boundary these are implemented on
$B=\operatorname{supp}(r)$ and extended arbitrarily to $A$, as licensed by Lemma~\ref{pt:lem:supprestrict}.  Padding record
alphabets if needed, Condition~\textup{(P5)} compares the two continuation profiles that differ only in branch $\bar x$.
If $F_r(\nu)>F_r(\nu')$, Condition~\textup{(P5)} gives $g(\nu)>g(\nu')$.  Conversely, if $g(\nu)\ge g(\nu')$ but
$F_r(\nu)<F_r(\nu')$, applying the strict part of Condition~\textup{(P5)} to the swapped branch comparison gives
$g(\nu')>g(\nu)$, a contradiction.  Hence $g(\nu)\ge g(\nu')$ iff $F_r(\nu)\ge F_r(\nu')$.

Since $|\operatorname{supp}(r)|\ge2$, local non-triviality on the support face and Lemmas~\ref{pt:lem:supprestrict} and~\ref{pt:lem:convnorm} imply $F_r$ is nonconstant.  The continuous affine functions $g$ and $F_r$ are therefore nonconstant
affine representatives of the same weak order on $\mathcal M_r$.  Herstein--Milnor uniqueness gives
\[
g(\nu)=\alpha\,F_r(\nu)+\gamma
\]
for some $\alpha>0$ and $\gamma\in\R$.

\emph{The coefficient $\alpha$ is independent of the fixed continuations.}  Write $g^{(1)}$ and $g^{(2)}$ for the
functionals arising from two different choices $\{H^x\}$ and $\{\widetilde H^x\}$ in the other branches.  Both
represent the same order as $F_r$, so $g^{(1)}=\alpha_1 F_r+\gamma_1$ and
$g^{(2)}=\alpha_2 F_r+\gamma_2$ with $\alpha_1,\alpha_2>0$.  To show $\alpha_1=\alpha_2$, consider the
difference $g^{(1)}-g^{(2)}=(\alpha_1-\alpha_2)F_r+(\gamma_1-\gamma_2)$.  For any
$\nu\in\mathcal M_r$,
\[
g^{(1)}(\nu)-g^{(2)}(\nu)
=
F_q\!\left(m(\bar x)\nu+\textstyle\sum_{x\ne\bar x}m(x)\mu_{r_x,H^x}\right)
-
F_q\!\left(m(\bar x)\nu+\textstyle\sum_{x\ne\bar x}m(x)\mu_{r_x,\widetilde H^x}\right).
\]
Both arguments are posterior laws in $\mathcal M_q$, differing only in the fixed non-$\bar x$ components.  Since $L_q$
has already been defined on the affine-hull difference space,
\[
g^{(1)}(\nu)-g^{(2)}(\nu)=L_q(\rho_1-\rho_2),
\]
where $\rho_i$ is the fixed non-$\bar x$ background law for $g^{(i)}$.  This difference is independent of $\nu$.  Since
$F_r$ is nonconstant in the present nondegenerate case, $\alpha_1=\alpha_2$.

\emph{Path-independence of the branch coefficient via tangent-space decomposition.}
We now prove that $\alpha/m(\bar x)$ depends only on the pair $(q,r_{\bar x})$, not on the first-stage experiment.
This is the key step that converts ordinal branchwise monotonicity into a cardinal aggregation formula.

\emph{Tangent space at a posterior.}  For a posterior $r\in\Delta(A)$ with support $B:=\operatorname{supp}(r)$, define
the \emph{tangent space}:
\[
V_r:=\operatorname{span}\{\nu-\nu':\nu,\nu'\in\mathcal M_r\}.
\]
Elements of $V_r$ are signed measures with total mass zero and barycentre zero.

\emph{Tangent-space identification.}  For any $r$ with full support, we claim
\[
V_r=W:=\{\eta\in\operatorname{span}\Delta_f(\Delta(A)):\eta(\Delta(A))=0,\ \textstyle\int p\,d\eta(p)=0\},
\]
the space of finite-support signed measures on $\Delta(A)$ with zero total mass and zero barycentre.
Indeed, if $\nu,\nu'\in\mathcal M_r$, then $\nu-\nu'$ has total mass $1-1=0$ and barycentre $r-r=0$, so $V_r\subset W$.
Conversely, given any $\eta\in W$: if $\eta=0$ then $\eta\in V_r$ trivially.  Otherwise, since $\eta$ has finite support,
	write its Jordan decomposition $\eta=\eta^+-\eta^-$ where $\eta^+,\eta^-\ge 0$ are finite-support positive measures
	with equal total mass $c>0$.  Normalise $\zeta^+:=\eta^+/c$ and $\zeta^-:=\eta^-/c$; both are finite-support probability
	measures on $\Delta(A)$ with the same barycentre
	$b:=\int p\,d\zeta^+(p)=\int p\,d\zeta^-(p)$.  Since $b\in\Delta(A)$, the set $\{a:b(a)>0\}$ is nonempty.  Because $r$
	has full support,
	\[
	\theta:=\min\Bigl\{1,\min_{a:b(a)>0}\frac{r(a)}{b(a)}\Bigr\}>0.
	\]
	Choose $0<\lambda<\theta$ and set
	\[
	\tau:=\frac{r-\lambda b}{1-\lambda}.
	\]
	Then $1-\lambda>0$, $\sum_a\tau(a)=1$, and $\tau(a)\ge0$ for every $a$: if $b(a)>0$ this follows from
	$\lambda<r(a)/b(a)$, while if $b(a)=0$ the numerator is $r(a)>0$.  Hence $\tau\in\Delta(A)$.  Define
	\[
	\nu:=\lambda\zeta^++(1-\lambda)\delta_\tau,
	\qquad
	\nu':=\lambda\zeta^-+(1-\lambda)\delta_\tau.
	\]
	Both are finite-support probability measures on $\Delta(A)$ with barycentre $\lambda b+(1-\lambda)\tau=r$, so
	$\nu,\nu'\in\mathcal M_r$.  Then $\nu-\nu'=\lambda(\zeta^+-\zeta^-)=(\lambda/c)\eta$, hence
	$\eta=(c/\lambda)(\nu-\nu')\in V_r$.  Thus $W\subset V_r$ and $V_r=W$.
Since $W$ depends only on the full-support property (not on the specific prior), $V_q=V_r=V_s=W$ for all full-support priors.

\emph{Case 1: Degenerate branch posterior ($|B|=1$).}  If $r=\delta_a$ for some action $a$, then
	$\mathcal M_{\delta_a}=\{\delta_{\delta_a}\}$ is a singleton: the only Bayes-plausible law with barycentre $\delta_a$
	is the point mass at $\delta_a$ itself.  Hence $V_r=\{0\}$, and every continuation experiment at prior $\delta_a$ gives
	the same posterior law $\delta_{\delta_a}$.  We set $F_{\delta_a}(\delta_{\delta_a})=0$.  For such branches, the term
	$m(x)\beta(q,r_x)F_{r_x}(\mu_{r_x,Q^x})$ is zero under any convention for $\beta(q,r_x)$.

For the aggregation formula below, assign any positive value to $\beta(q,\delta_a)$; the choice is not behaviourally
identified because the corresponding continuation value is always zero.


\emph{Case 2: Non-full-support branch posterior ($1<|B|<|A|$).}  If $r$ has support $B\subsetneq A$ with $|B|\ge 2$,
apply Lemma~\ref{pt:lem:supprestrict}.  Continuation comparisons at prior $r$ reduce to those at $r|_B\in\Delta(B)$,
with representative $F_r:=F_{r|_B}^B\circ\pi_B$.

For any finite action set $C$, write
\[
W_C:=\{\sigma\in\operatorname{span}\Delta_f(\Delta(C)):\sigma(\Delta(C))=0,\ \textstyle\int_{\Delta(C)}p\,d\sigma(p)=0\}.
\]
Let $\jmath_B:\Delta(B)\to\Delta(A)$ be the face inclusion, extending each posterior by zero outside $B$, and let
$i_B:=(\jmath_B)_\#$ act on finite signed posterior laws.  By Lemma~\ref{pt:lem:supprestrict}(iv),
$F_r^A=F_{r|_B}^B\circ\pi_B$, so on face directions
\[
L_r^A(i_B\sigma)=L_{r|_B}^B(\sigma)\qquad(\sigma\in W_B).
\]
Since $|B|\ge2$ and $r|_B$ is full support on $B$, local non-triviality and Lemma~\ref{pt:lem:convnorm} give
\[
L_{r|_B}^B(\chi_{r|_B}-\delta_{r|_B})>0,
\]
so $L_{r|_B}^B\ne0$.

\emph{Full-to-face scalar.}  Fix any full-support $q\in\Delta(A)$.  The boundary posterior $r$ is reachable from $q$:
choose
\[
0<\varepsilon<\min\Bigl\{1,\min_{b\in B}\frac{q(b)}{r(b)}\Bigr\}
\]
and set $P_1(1\mid a)=\varepsilon r(a)/q(a)$, with $r(a)=0$ for $a\notin B$.  Then branch $1$ has positive probability
$\varepsilon$ and posterior $r$.  For any nonzero $\sigma\in W_B$, the tangent-space identification on the face $B$ gives
$\sigma=t(\nu-\nu')$ for some $t>0$ and $\nu,\nu'\in\mathcal M_{r|_B}^B$.  Implement $\nu,\nu'$ by continuation channels
on $B$ and extend them arbitrarily to $A$.  Lemma~\ref{pt:lem:supprestrict} identifies their comparison at the boundary prior
$r$ with their comparison at $r|_B$.

Condition~\textup{(P5)} applied to the branch reaching $r$, with all other branches fixed, gives the positive sign direction
for feasible differences.  The negative direction follows by swapping the two continuation laws, and equality follows by
applying the weak part of Condition~\textup{(P5)} in both directions.  Extending by positive scalar multiples from feasible
differences to all of $W_B$, $L_q^A\circ i_B$ and $L_{r|_B}^B$ have the same sign partition on $W_B$.  Since
$L_{r|_B}^B\ne0$, both are nonzero and there is a unique $\beta_A(q,r)>0$ such that
\[
L_q^A(i_B\sigma)=\beta_A(q,r)L_{r|_B}^B(\sigma)\qquad(\sigma\in W_B).
\]
This is the ambient full-to-face coefficient used for boundary branches.

\emph{Case 3: Full-support branch posterior ($B=A$).}  This is the main case.  Suppose $r$ has full support on $A$.

\emph{Claim 1: aggregate change from branch variation.}  If a first-stage branch reaches posterior $r$ with probability
$m$, and the branch continuation law changes from $\nu$ to $\nu'$ (both in $\mathcal M_r$), then the aggregate posterior
law changes by $m(\nu-\nu')\in m\cdot V_r$.  This difference depends only on $(m,\nu,\nu',r)$---not on which experiment
produced the branch, nor on the other branches.

\emph{Proof of Claim 1.}  The aggregate changes from $m\nu+\rho$ to $m\nu'+\rho$, where $\rho$ collects the other
branches.  The difference is $m(\nu-\nu')$.

\emph{Claim 2: $L_q|_{V_r}$ is a positive scalar multiple of $L_r|_{V_r}$.}

\emph{Proof of Claim 2.}  We first show that every direction in $V_r$ is proportional to a feasible difference.
By the Tangent-space identification above, any nonzero $\sigma\in V_r$ can be written as
$\sigma=t(\nu-\nu')$ for some $t>0$ and $\nu,\nu'\in\mathcal M_r$.  Hence every nonzero $\sigma\in V_r$ is a positive
scalar multiple of some feasible difference.

By Condition~\textup{(P5)} (branchwise monotonicity), if $\nu\succ_r\nu'$ (i.e.\ $F_r(\nu)>F_r(\nu')$, equivalently
$L_r(\nu-\nu')>0$), then the aggregate comparison ranks the $\nu$-continuation strictly higher.  Since the aggregate
value difference is $L_q(m(\nu-\nu'))=mL_q(\nu-\nu')$ with $m>0$, we have $L_q(\nu-\nu')>0$ whenever $L_r(\nu-\nu')>0$.

\emph{Reverse and equality directions.}  If $L_r(\nu-\nu')<0$, completeness of $\succeq_r$ gives $\nu'\succ_r\nu$, and
Condition~\textup{(P5)} applied to the reversed comparison gives $L_q(\nu'-\nu)>0$, i.e.\ $L_q(\nu-\nu')<0$.  If $L_r(\nu-\nu')=0$, then
$\nu\sim_r\nu'$, so both $\nu\succeq_r\nu'$ and $\nu'\succeq_r\nu$; applying Condition~\textup{(P5)} weakly in both directions gives
$L_q(\nu-\nu')\ge 0$ and $L_q(\nu-\nu')\le 0$, hence $L_q(\nu-\nu')=0$.

Since every nonzero $\sigma\in V_r$ is proportional to some feasible difference, the sign agreement extends from
feasible differences to all of $V_r$: for any $\sigma\in V_r$, $L_q(\sigma)>0$ iff $L_r(\sigma)>0$, $L_q(\sigma)<0$
iff $L_r(\sigma)<0$, and $L_q(\sigma)=0$ iff $L_r(\sigma)=0$.

Two nonzero linear functionals on a vector space that have the same positive half-space must be positive scalar
multiples of each other.  (If $L_r\ne 0$, then $\ker L_q=\ker L_r$ and they lie in the same half-space, so $L_q=cL_r$
for some $c>0$.)  Since $r$ has full support and local non-triviality gives $F_r(\chi_r)>F_r(\delta_r)=0$, we have
$L_r\ne 0$.  Hence there exists $\beta(q,r)>0$ such that
\[
L_q(\sigma)=\beta(q,r)\,L_r(\sigma)
\qquad\text{for all }\sigma\in V_r.
\]

\emph{Path-independence.}  Consider any first-stage experiment reaching posterior $r$ in some branch with probability
$m$.  Changing the branch continuation from $\nu$ to $\nu'$ changes the aggregate value by
\[
L_q(m(\nu-\nu'))=m\,\beta(q,r)\,L_r(\nu-\nu')=m\,\beta(q,r)\,(F_r(\nu)-F_r(\nu')).
\]
The coefficient $m\,\beta(q,r)$ depends only on $(q,r,m)$.  In particular, $\beta(q,r)$ is independent of which
experiment reached $r$ and independent of the other branches.

Comparing with the earlier notation $g(\nu)=\alpha F_r(\nu)+\gamma$: the value difference is
$g(\nu)-g(\nu')=\alpha(F_r(\nu)-F_r(\nu'))$, so $\alpha=m\,\beta(q,r)$ and hence $\alpha/m=\beta(q,r)$ universally.
The coefficient $\beta(q,r_{\bar x})$ is defined as the scalar from $L_q|_{V_{r_{\bar x}}}=\beta(q,r_{\bar x})L_{r_{\bar x}}|_{V_{r_{\bar x}}}$;
it depends only on $q$ and $r_{\bar x}$, not on the experiment that produced the branch.

\emph{Aggregation formula.}  Let
\[
\mu^0:=\sum_{x:m(x)>0}m(x)\delta_{r_x}=\mu_{q,P_1}.
\]
For each positive-probability branch let $B_x:=\operatorname{supp}(r_x)$.  If $B_x=A$, set
$v_x:=\mu_{r_x,Q^x}-\delta_{r_x}\in W_A$.  If $1<|B_x|<|A|$, set
\[
\bar v_x:=\pi_{B_x}(\mu_{r_x,Q^x})-\delta_{r_x|_{B_x}}\in W_{B_x},
\qquad
v_x:=i_{B_x}\bar v_x\in W_A.
\]
If $|B_x|=1$, set $v_x:=0$.  Then
\[
\mu_{q,P_1\triangleright\{Q^x\}}=\mu^0+\sum_{x:m(x)>0}m(x)v_x.
\]
Since $V_q=W_A$, affinity of $F_q$ gives
\[
F_q(\mu_{q,P_1\triangleright\{Q^x\}})-F_q(\mu^0)
=\sum_{x:m(x)>0}m(x)L_q^A(v_x).
\]
For a full-support branch $B_x=A$, Case~3 gives
\[
L_q^A(v_x)=\beta(q,r_x)L_{r_x}^A(v_x)
=\beta(q,r_x)F_{r_x}(\mu_{r_x,Q^x}).
\]
For a nondegenerate boundary branch $1<|B_x|<|A|$, Case~2 gives
\[
L_q^A(v_x)=L_q^A(i_{B_x}\bar v_x)
=\beta_A(q,r_x)L_{r_x|_{B_x}}^{B_x}(\bar v_x)
=\beta_A(q,r_x)F_{r_x}(\mu_{r_x,Q^x}).
\]
The last equalities use zero normalisation: Lemma~\ref{pt:lem:convnorm} in full support, Lemma~\ref{pt:lem:supprestrict} plus
Lemma~\ref{pt:lem:convnorm} on the support face in the boundary case, and the singleton convention when $|B_x|=1$.  Therefore
\[
F_q(\mu_{q,P_1\triangleright\{Q^x\}})
=
F_q(\mu_{q,P_1})+\sum_{x:m(x)>0}m(x)\,\beta(q,r_x)\,F_{r_x}(\mu_{r_x,Q^x}),
\]
where $\beta(q,r_x)$ denotes the Case~3 coefficient when $B_x=A$, the Case~2 full-to-face coefficient when
$1<|B_x|<|A|$, and any positive singleton convention when $|B_x|=1$.
\end{proof}

\begin{lemma}[Branch-coefficient cocycle and normalised chain rule]\label{pt:lem:chain}
Let the full-support and full-to-face coefficients be those constructed in Lemma~\ref{pt:lem:branchagg}, applied in the
relevant ambient simplex or intrinsically on the relevant support face.  If $q$ has full support and $r,s$ are
nondegenerate priors with
\[
\operatorname{supp}(s)\subseteq\operatorname{supp}(r)\subseteq\operatorname{supp}(q),
\]
then the coefficients satisfy the cocycle identity
\[
\beta(q,s)=\beta(q,r)\beta(r,s).
\]
Hence $\beta(q,r)=a_q/a_r$ for positive scales $a_q$, with singleton scales fixed by convention.  The rescaled values
\[
\widehat F_q(\mu):=\frac{F_q(\mu)}{a_q}
\]
satisfy the exact chain rule: for every first-stage channel $P_1$ at full-support $q$ and every continuation profile
$\{Q^x\}_{x:m(x)>0}$,
\[
\widehat F_q(\mu_{q,P_1\triangleright\{Q^x\}})
=
\widehat F_q(\mu_{q,P_1})
\;+\;
\sum_{x:m(x)>0}m(x)\,\widehat F_{r_x}(\mu_{r_x,Q^x}).
\]
\end{lemma}

\begin{proof}

\emph{Full-support cocycle.}  First consider reachable full-support $q,r,s\in\operatorname{int}\Delta(A)$ with $|A|\ge2$.
For such priors, all three tangent spaces coincide.
By the tangent-space identification in Lemma~\ref{pt:lem:branchagg}, $V_q=V_r=V_s=W$ for all full-support priors.
The full-support coefficient identities constructed there share a common domain $W$:
\[
L_q|_W=\beta(q,r)L_r|_W,
\qquad
L_r|_W=\beta(r,s)L_s|_W.
\]
Substituting the second into the first:
\[
L_q|_W=\beta(q,r)\beta(r,s)L_s|_W.
\]
Comparing with $L_q|_W=\beta(q,s)L_s|_W$ and using $L_s\ne 0$:
\[
\beta(q,s)=\beta(q,r)\beta(r,s).
\]

\emph{Reachability.}  The above uses the full-support coefficient construction for the pairs $(q,r)$, $(r,s)$, and
$(q,s)$.  Each invocation requires the pair to be reachable: the second prior must be a possible posterior of some
experiment at the first prior.
For any full-support $q,r$, construct a two-record experiment $P_1$ with $P_1(1\mid a)=\varepsilon r(a)/q(a)$ for
$0<\varepsilon<\min\{1,\min_a q(a)/r(a)\}$.  Record~1 has probability $\varepsilon$ and posterior~$r$; record~0 has posterior
$(q-\varepsilon r)/(1-\varepsilon)$.  Hence any full-support $r$ is reachable from any full-support $q$, so the construction
applies to all full-support pairs.  The cocycle therefore holds for all full-support triples.

Fix a full-support basepoint $q_0$, set $a_{q_0}=1$, and for full-support $q$ define
\[
a_q:=\frac{1}{\beta(q_0,q)}.
\]
By the full-support cocycle, this is equivalently $a_q=\beta(q,q_0)$, and
\[
\beta(q,r)=\frac{\beta(q,q_0)}{\beta(r,q_0)}=\frac{a_q}{a_r}.
\]

\emph{Boundary coefficients and face scales.}  If $r$ is nondegenerate with support $B\subsetneq A$, Lemma~\ref{pt:lem:branchagg}
constructs the ambient full-to-face coefficient $\beta_A(q,r)>0$ by
\[
L_q^A i_B=\beta_A(q,r)L_{r|_B}^B.
\]
When the ambient action set is clear, write this coefficient as $\beta(q,r)$.

When nondegenerate boundary priors $r,s$ have common support $B$, $\beta_A(r,s)$ means the unique $c>0$ satisfying
\[
L_r^A(i_B\sigma)=c\,L_s^A(i_B\sigma)\qquad(\sigma\in W_B).
\]
This is a scalar on the common face, not on all of $W_A$.  Since
\[
L_r^A\circ i_B=L_{r|_B}^B,\qquad L_s^A\circ i_B=L_{s|_B}^B,
\]
and the intrinsic face scalar satisfies
$L_{r|_B}^B(\sigma)=\beta_B(r|_B,s|_B)L_{s|_B}^B(\sigma)$, we have
\[
\beta_A(r,s)=\beta_B(r|_B,s|_B).
\]

Fix the same full-support basepoint $q_0$ on $\Delta(A)$.  If $r_0$ is any nondegenerate face basepoint with support $B$,
then for every nondegenerate $r$ with support $B$,
\[
\beta_A(q_0,r)=\beta_A(q_0,r_0)\,\beta_B(r_0|_B,r|_B).
\]
Indeed, both sides are the coefficient multiplying $L_{r|_B}^B$ to obtain $L_{q_0}^A\circ i_B$, and
$L_{r|_B}^B\ne0$.  Define the ambient boundary scale by
\[
a_r:=\frac{1}{\beta_A(q_0,r)}.
\]
If the face scale is based at $r_0$ with $a_{r_0}:=1/\beta_A(q_0,r_0)$, then
\[
a_{r|_B}^B:=\frac{a_{r_0}}{\beta_B(r_0|_B,r|_B)}
=\frac{1}{\beta_A(q_0,r_0)\beta_B(r_0|_B,r|_B)}
=\frac{1}{\beta_A(q_0,r)}
=a_r.
\]
Thus the ambient and intrinsic scales agree: treating $r$ as a boundary posterior in $\Delta(A)$ or as a full-support prior
on $\Delta(B)$ yields the same normalised representative $\widehat F_r$.

\emph{Extension to boundary posteriors.}  If $r$ is nondegenerate with support $B\subsetneq A$, Lemma~\ref{pt:lem:branchagg}
defines
$\beta(q,r)=\beta_A(q,r)$ by
\[
L_q^A i_B=\beta_A(q,r)L_{r|_B}^B.
\]
Using the full-support identity $L_q^A=\beta(q,q_0)L_{q_0}^A$ and restricting to $i_B(W_B)$,
\[
L_q^A i_B=\beta(q,q_0)L_{q_0}^A i_B
=\beta(q,q_0)\beta_A(q_0,r)L_{r|_B}^B.
\]
Comparing with the definition of $\beta_A(q,r)$ and using $L_{r|_B}^B\ne0$ gives
\[
\beta(q,r)=\beta(q,q_0)\beta(q_0,r)=\frac{a_q}{a_r},
\]
where $a_r=1/\beta(q_0,r)$ by the boundary-scale definition above.

\emph{Singleton scale conventions.}  If $r=\delta_a$ is a singleton prior, set $a_{\delta_a}:=a_{q_0}=1$ and
$\widehat F_{\delta_a}\equiv0$.  Define $\beta(q,\delta_a):=a_q/a_{\delta_a}$ as a convention.  No behavioural branch
coefficient is identified at a singleton target, but singleton terms in the aggregation and chain formulas are always zero.

\emph{Mixed-support cocycle.}  Let $r,s$ be nondegenerate with nested supports
$C:=\operatorname{supp}(s)\subseteq B:=\operatorname{supp}(r)$, and write $s|_B$ for $s$ viewed as a boundary prior on
the action set $B$.  Interpret
\[
\beta(r,s):=\beta_B(r|_B,s|_B),
\]
the intrinsic full-to-face coefficient on action set $B$.  Let $i_C^A:W_C\to W_A$ and $i_C^B:W_C\to W_B$ be the face
inclusions.  For every $\sigma\in W_C$,
\[
L_q^A i_C^A\sigma
=\beta(q,r)L_{r|_B}^B i_C^B\sigma
=\beta(q,r)\beta(r,s)L_{s|_C}^C\sigma.
\]
The direct full-to-face definition from Lemma~\ref{pt:lem:branchagg} for the pair $(q,s)$ gives
\[
L_q^A i_C^A\sigma=\beta(q,s)L_{s|_C}^C\sigma.
\]
Since $L_{s|_C}^C\ne0$, $\beta(q,s)=\beta(q,r)\beta(r,s)$.  If
$\operatorname{supp}(s)\not\subseteq\operatorname{supp}(r)$, then $\beta(r,s)$ is not behaviourally defined; singleton
targets are also excluded from the behavioural cocycle because their scales are conventions only.

\emph{Chain rule.}  Applying Lemma~\ref{pt:lem:branchagg} and substituting $\beta(q,r)=a_q/a_r$ gives
\[
F_q(\mu_{q,P_1\triangleright\{Q^x\}})
=
F_q(\mu_{q,P_1})+\sum_{x:m(x)>0}m(x)\,\frac{a_q}{a_{r_x}}\,F_{r_x}(\mu_{r_x,Q^x}).
\]
Dividing by $a_q$ and writing $\widehat F_q(\mu):=F_q(\mu)/a_q$ yields
\[
\widehat F_q(\mu_{q,P_1\triangleright\{Q^x\}})
=
\widehat F_q(\mu_{q,P_1})+\sum_{x:m(x)>0}m(x)\,\widehat F_{r_x}(\mu_{r_x,Q^x}).
\qedhere
\]
\end{proof}

\begin{lemma}[Coherent relabelling and face scales]\label{pt:lem:facescales}
The scales produced by Lemma~\ref{pt:lem:chain} can be chosen simultaneously across all finite action sets so that:
\begin{enumerate}
\item for every bijection $\pi:A\to A'$,
\[
a^{A'}_{q\pi}=a^A_q;
\]
\item if $\iota:B\hookrightarrow A$ is an injection with $|B|\ge2$, $r$ has full support on $B$, and $q$ has full support
on $A$, then the full-to-face branch coefficient for reaching the face $\iota(B)$ is
\begin{equation}
\beta_\iota(q,r)=\frac{a^A_q}{a^B_r}.
\tag{FS}\label{pt:eq:facescale}
\end{equation}
\end{enumerate}
Singleton target scales remain a convention because every continuation value on a singleton fibre is zero.
\end{lemma}

\begin{proof}
For every nondegenerate finite action set $A$, first choose the full-support scales supplied by the cocycle part of
Lemma~\ref{pt:lem:chain}, so that
\[
\beta_A(q,q')=\frac{a^A_q}{a^A_{q'}}
\qquad(q,q'\in\operatorname{int}\Delta(A)).
\]
They are unique up to multiplication by one positive constant depending on $A$.  Normalise initially by
$a^A_{u_A}=1$, where $u_A$ is uniform.  Exact relabelling invariance of the representatives implies relabelling
invariance of the uniquely determined branch coefficients.  Since bijections carry uniform priors to uniform priors,
the initial scales are exactly invariant under bijections.

Let $\iota:B\hookrightarrow A$ be an injection between nondegenerate finite sets.  By relabelling $B$ onto its image
$\iota(B)$ and using Corollary~\ref{pt:cor:permutationinvariance}, it is enough to define the coefficient for the inclusion
of the face $\iota(B)$; the notation below keeps the original symbol $\iota$.  Let $i_\iota$ denote the induced map on
signed posterior-law tangent spaces.  For full-support $q\in\Delta(A)$ and $r\in\Delta(B)$, let
$\beta_\iota(q,r)>0$ be the full-to-face scalar characterised by
\[
L_q^A\circ i_\iota=\beta_\iota(q,r)L_r^B
\quad\hbox{on }W_B.
\]
The number
\begin{equation}
K(\iota):=\beta_\iota(q,r)\frac{a^B_r}{a^A_q}
\tag{K}\label{pt:eq:embeddingdefect}
\end{equation}
is independent of both $q$ and $r$.  Independence of $q$ follows by comparing two full-support ambient priors and using
$L_q^A=(a_q^A/a_{q'}^A)L_{q'}^A$ on $W_A$.  Independence of $r$ follows from same-face compatibility: if $r,s$ have
full support on $B$, then $L_r^B=(a_r^B/a_s^B)L_s^B$, and hence
$\beta_\iota(q,s)=\beta_\iota(q,r)a_r^B/a_s^B$.

The defects multiply under composition.  If
$C\xhookrightarrow{\jmath}B\xhookrightarrow{\iota}A$ and all three sets in the comparison are nondegenerate, then
restricting first to $B$ and then to $C$ gives, for full-support $q,r,s$ on $A,B,C$ respectively,
\[
\beta_{\iota\circ\jmath}(q,s)=\beta_\iota(q,r)\beta_\jmath(r,s).
\]
Substituting the definition of $K$ and cancelling the intermediate scale $a_r^B$ gives
\begin{equation}
K(\iota\circ\jmath)=K(\iota)K(\jmath).
\tag{KC}\label{pt:eq:embeddingcocycle}
\end{equation}
Exact relabelling invariance shows that $K(\iota)$ depends only on $n:=|A|$ and $m:=|B|$; write it as $K_{n,m}$.
Bijections give $K_{n,n}=1$, and
\[
K_{n,\ell}=K_{n,m}K_{m,\ell}
\qquad(2\le\ell\le m\le n).
\]
Put $t_n:=K_{n,2}$, with $t_2=1$.  Then $K_{n,m}=t_n/t_m$.  Finally replace every scale on an $n$-point action set by
\[
\widetilde a_q^A:=t_n a_q^A.
\]
Under this change, the defect in~\eqref{pt:eq:embeddingdefect} becomes
\[
\widetilde K_{n,m}=K_{n,m}\frac{t_m}{t_n}=1.
\]
The rescaling is common to all priors on a fixed action set, so it leaves all within-set cocycle ratios unchanged; because
$t_n$ depends only on cardinality, it also preserves exact relabelling invariance.  Dropping tildes gives
\eqref{pt:eq:facescale}.  Singleton target scales may be set arbitrarily because all singleton continuation terms are zero.
\end{proof}

\paragraph{Stage 5: compatibility collapse and entropy identification.}
\addcontentsline{toc}{subsubsection}{Stage 5: compatibility collapse and entropy identification}
The product and sequential evaluations of full revelation must agree.  Their
compatibility produces a positive grouping-weight equation; evaluating one
distribution by two groupings forces the product interaction to vanish and
the chain scale to be universal.  Entropy reduction and Faddeev's recursion
then identify the remaining functional.

\begin{lemma}[Interaction collapse and universal chain scale]\label{pt:lem:scalecoherence}
The universal interaction coefficient in~\eqref{pt:eq:QA} is zero.  Moreover, there exists $a>0$ such that the coherently
chosen chain scales satisfy $a_q=a$ for every nondegenerate full-support prior on every finite action set.  The same
value may be assigned to singleton priors.  Consequently,
\begin{equation}
F_{p\otimes r}(\mu\otimes\nu)=F_p(\mu)+F_r(\nu)
\tag{PA}\label{pt:eq:PArecovered}
\end{equation}
and the branch formula has a universal chain scale.
\end{lemma}

\begin{proof}
Define, for every finite prior $q$ read on its support,
\[
H(q):=F_q(\chi_q),
\qquad
Z(q):=1+\kappa H(q).
\]
For a nondegenerate prior, local non-triviality gives $H(q)>0$, and~\eqref{pt:eq:QAslope} gives
\begin{equation}
Z(q)>0.
\tag{Z+}\label{pt:eq:Zpositive}
\end{equation}
For a singleton, $H(q)=0$ and $Z(q)=1$.

\emph{Step 1: product revelation links the chain scale to $Z$.}
Let $q\in\Delta(A)$ and $r\in\Delta(B)$ be nondegenerate and full support.  Quasi-additivity applied to full revelation
gives
\begin{equation}
H(q\otimes r)=H(q)+H(r)+\kappa H(q)H(r).
\tag{QH}\label{pt:eq:QH}
\end{equation}
Compute the same value sequentially.  First reveal the $A$-coordinate and then reveal the $B$-coordinate in every branch.
The first-stage value is $F_{q\otimes r}(\chi_q\otimes\delta_r)=H(q)$ by~\eqref{pt:eq:QA}.  After observing $a$, the
posterior is $e_a\otimes r$ on the face $\{a\}\times B$; support restriction and relabelling identify its continuation
full-revelation value with $H(r)$.  For the injection
\[
\iota_a:B\hookrightarrow A\times B,\qquad b\mapsto(a,b),
\]
Lemma~\ref{pt:lem:facescales} gives the branch coefficient
\[
\beta_{\iota_a}(q\otimes r,r)=\frac{a^{A\times B}_{q\otimes r}}{a^B_r}.
\]
The branch weights are $q(a)$ and the same coefficient and continuation value appear in every branch, so
$\sum_a q(a)=1$.  Lemma~\ref{pt:lem:branchagg} therefore yields
\begin{equation}
H(q\otimes r)=H(q)+\frac{a_{q\otimes r}}{a_r}H(r).
\tag{SA}\label{pt:eq:seqA}
\end{equation}
Comparing~\eqref{pt:eq:QH} and~\eqref{pt:eq:seqA}, and using $H(r)>0$,
\begin{equation}
\frac{a_{q\otimes r}}{a_r}=Z(q).
\tag{R1}\label{pt:eq:ratio1}
\end{equation}
Revealing $B$ first gives symmetrically
\begin{equation}
\frac{a_{q\otimes r}}{a_q}=Z(r).
\tag{R2}\label{pt:eq:ratio2}
\end{equation}
Hence $a_rZ(q)=a_qZ(r)$, so there is a universal constant $C>0$ such that
\begin{equation}
a_q=CZ(q)
\tag{AS}\label{pt:eq:ascaleZ}
\end{equation}
for every nondegenerate full-support prior on every finite action set.  Assign $a_q:=C$ at singleton priors, where
$Z(q)=1$ and all continuation values are zero.

If $\kappa=0$, then $Z\equiv1$, equation~\eqref{pt:eq:ascaleZ} already gives universal scales, and~\eqref{pt:eq:QA} is exact
product additivity.  For the remaining argument suppose $\kappa\ne0$.

\emph{Step 2: the induced weight equation.}
Let $q$ be full support on $A$, let $\mathcal P=\{A_1,\ldots,A_m\}$ be a partition into nonempty cells, and put
\[
p_i:=q(A_i),
\qquad
q^i:=q(\cdot\mid A_i),
\qquad
p:=(p_1,\ldots,p_m).
\]
Let $G_{\mathcal P}$ reveal the block.  Undetectable-distinctions neutrality gives
$(q,G_{\mathcal P})\sim(p,\Id_{\{1,\ldots,m\}})$.  The early cross-prior representation
Lemma~\ref{pt:lem:blockbridge}, importantly for the \emph{unscaled} $F$'s, therefore gives
\begin{equation}
F_q(\mu_{q,G_{\mathcal P}})=H(p).
\tag{BG}\label{pt:eq:blockvalue}
\end{equation}
Append full revelation in every branch.  The compound posterior law is $\chi_q$, so the chain formula and
\eqref{pt:eq:ascaleZ} give
\begin{equation}
H(q)=H(p)+Z(q)\sum_{i=1}^m p_i\frac{H(q^i)}{Z(q^i)}.
\tag{GR}\label{pt:eq:groupingH}
\end{equation}
Define
\[
w(q):=\frac1{Z(q)}>0.
\]
Using $\kappa H(x)=Z(x)-1$ in~\eqref{pt:eq:groupingH} gives
\[
\frac{Z(p)}{Z(q)}=\sum_{i=1}^m p_i\frac1{Z(q^i)}.
\]
Equivalently,
\begin{equation}
w(q)=w(p)\sum_{i=1}^m p_iw(q^i).
\tag{W}\label{pt:eq:weight}
\end{equation}
No continuity assertion is needed here.

\emph{Step 3: a two-grouping argument forces $w\equiv1$.}
All distributions in this paragraph are read on their positive-probability supports, so the weight equation applies after
support restriction.  First,~\eqref{pt:eq:weight} applied to a product distribution gives
\begin{equation}
w(u\otimes v)=w(u)w(v),
\tag{M}\label{pt:eq:wmult}
\end{equation}
because the first-coordinate block probabilities are $u$ and every conditional distribution is $v$.

Take arbitrary finite probability vectors $u$ and $v$, possibly on different alphabets, and write
\[
x:=w(u),\qquad y:=w(v),\qquad p:=(1/2,1/2).
\]
On the labelled disjoint union $(U\times U)^0\sqcup(V\times V)^1$, form the distribution
\[
T:=\tfrac12(u\otimes u)^0\;\sqcup\;\tfrac12(v\otimes v)^1.
\]
Grouping first by the top block and using~\eqref{pt:eq:wmult},
\begin{equation}
w(T)=w(p)\frac{x^2+y^2}{2}.
\tag{E1}\label{pt:eq:twoeval1}
\end{equation}
Now group first by the pair consisting of the top-block label and the first within-block coordinate.  Its marginal is
\[
S:=\tfrac12u^0\;\sqcup\;\tfrac12v^1,
\]
and~\eqref{pt:eq:weight} gives
\[
w(S)=w(p)\frac{x+y}{2}.
\]
Conditional on a point in the first part of $S$, the remaining coordinate has distribution $u$; conditional on a point in
the second part, it has distribution $v$.  A second application of~\eqref{pt:eq:weight} therefore gives
\begin{equation}
w(T)=w(S)\frac{x+y}{2}
=w(p)\left(\frac{x+y}{2}\right)^2.
\tag{E2}\label{pt:eq:twoeval2}
\end{equation}
Since $w(p)>0$, comparison of~\eqref{pt:eq:twoeval1} and~\eqref{pt:eq:twoeval2} yields
\[
\frac{x^2+y^2}{2}=\left(\frac{x+y}{2}\right)^2,
\]
so $(x-y)^2=0$.  As $u$ and $v$ were arbitrary, $w$ is constant on all finite probability vectors.  At a singleton prior,
$H=0$, hence $w=1$.  Therefore $w(q)\equiv1$.

Thus $Z(q)=1$ and $\kappa H(q)=0$ for every $q$.  Local non-triviality gives a nondegenerate $q$ with $H(q)>0$, hence
$\kappa=0$, contradicting the temporary alternative unless the interaction coefficient is zero.  Equation
\eqref{pt:eq:ascaleZ} now gives $a_q=C$ universally.  Finally~\eqref{pt:eq:QA} reduces to exact product additivity
\eqref{pt:eq:PArecovered}.
\end{proof}

\begin{lemma}[Entropy reduction and Faddeev identification]\label{pt:lem:faddeevsketch}
By Lemma~\ref{pt:lem:scalecoherence}, $a_q\equiv a$ is universal.  Define
\[
\widehat F_q:=F_q/a,
\qquad
H(q):=\widehat F_q(\chi_q),
\]
the rescaled value of full revelation.  Then, for every channel $P$,
\[
\widehat F_q(\mu_{q,P})
=
H(q)-\sum_{x:m_{q,P}(x)>0}m_{q,P}(x)\,H(r_x^{\,q,P}).
\tag{ER}
\]
Moreover $H$ satisfies Faddeev's recursion, so
\[
H(q)=\alpha \Sh(q)
\]
for some $\alpha>0$.  Consequently,
\[
\widehat F_q(\mu_{q,P})=\alpha I_{q,P}(A;X).
\]
\end{lemma}

\begin{proof}
\emph{Entropy-reduction formula.}  Boundary values of $H$ are read through support restriction: if $q$ has support
$B$, then $H(q):=H(q|_B)$, and $H(\delta_a)=0$.  Let $P:A\to\Delta(X)$ be any channel at full-support prior $q$.  Append
full revelation $\Id_A:A\to\Delta(A)$ after $P$: the compound $P\triangleright\{\Id_A\}_x$ (using $\Id_A$ in every branch)
has $\mu_{r_x,\Id_A}=\chi_{r_x}$ and the overall posterior law is full revelation $\chi_q$.  Applying the chain rule
(Lemma~\ref{pt:lem:chain}) gives
\[
\widehat F_q(\chi_q)
=
\widehat F_q(\mu_{q,P})+\sum_{x:m(x)>0}m(x)\,\widehat F_{r_x}(\chi_{r_x}).
\]
Hence
\[
\widehat F_q(\mu_{q,P})=H(q)-\sum_{x:m(x)>0}m(x)\,H(r_x).
\]
If $q$ is on the boundary, restrict $q$ and $P$ to $B=\operatorname{supp}(q)$, apply the full-support formula on
$\Delta(B)$, and pull the value and comparison back by Lemma~\ref{pt:lem:supprestrict}; zero-probability rows and records
do not affect either side of the displayed identity.

\emph{Scaled cross-prior bridge.}
Lemma~\ref{pt:lem:blockbridge} was proved for the unscaled values $F_q$.  Since Lemma~\ref{pt:lem:scalecoherence} has made
$a_q\equiv a$ universal, it is equivalently valid for
\[
V(q,P):=\widehat F_q(\mu_{q,P})=F_q(\mu_{q,P})/a:
\]
for all finite priors and channels,
\[
q^0\succeq_{P\sqcup Q}r^1
\quad\Longleftrightarrow\quad
V(q,P)\ge V(r,Q).
\]

\emph{Faddeev recursion.}  First let $q$ be full support and let $\mathcal P=\{A_1,\dots,A_m\}$ be a partition into
nonempty cells.  Then $p_i:=q(A_i)>0$ for every $i$.  Let $G_{\mathcal P}:A\to\Delta(\{1,\dots,m\})$ reveal only the
block index: $G_{\mathcal P}(i\mid a)=\1_{a\in A_i}$.  Write
\[
p_i:=q(A_i),\qquad q^i:=q(\cdot\mid A_i),
\]
and let $p=(p_1,\dots,p_m)\in\Delta(\{1,\dots,m\})$.  Consider the deterministic kernel
$S:A\to\Delta(\{1,\dots,m\})$ given by
$S(i\mid a)=\1_{a\in A_i}$, so $qS=p$.

By Lemma~\ref{pt:lem:neutrality}, since $G_{\mathcal P}$ is $S$-measurable,
\[
(q,G_{\mathcal P})\sim(p,\Id_{\{1,\dots,m\}})
\]
in the two-block comparison.  Applying the scaled form of Lemma~\ref{pt:lem:blockbridge} to the forward comparison gives
$V(q,G_{\mathcal P})\ge V(p,\Id_{\{1,\dots,m\}})$.  Applying the reverse-block form from Lemma~\ref{pt:lem:blockcoh} and then
the same bridge to the reverse comparison gives $V(p,\Id_{\{1,\dots,m\}})\ge V(q,G_{\mathcal P})$.  Therefore
\[
\widehat F_q(\mu_{q,G_{\mathcal P}})
=\widehat F_p(\chi_p)=H(p).
\]
The entropy-reduction formula applied to $G_{\mathcal P}$ gives
\[
H(q)=\widehat F_q(\chi_q)
=\widehat F_q(\mu_{q,G_{\mathcal P}})+\sum_{i=1}^m p_i\widehat F_{q^i}(\chi_{q^i})
=H(p)+\sum_{i=1}^m p_iH(q^i).
\]
This is Faddeev's recursion.  For boundary $q$, restrict to $\operatorname{supp}(q)$, apply the same full-support recursion
to the nonempty positive-mass intersections, and pull back by Lemma~\ref{pt:lem:supprestrict}; equivalently, the recursion is
read over positive-mass blocks only.

We now verify that $H$ is continuous on each $\Delta_n:=\{q\in\Delta(A):|A|=n\}$.

\emph{Binary continuity via erasure calibrators.}
Define the binary entropy function $h:[0,1]\to\R$ by $h(t):=H(t,1-t)$.  We show $h$ is continuous using
prior-independent calibrators.  By Lemma~\ref{pt:lem:convnorm}, support restriction, and $a>0$, $H(q)\ge0$ for every $q$,
so $h(t)\ge0$ on $[0,1]$.  If $c<0$, then
\[
\{t:h(t)\ge c\}=[0,1],
\qquad
\{t:h(t)\le c\}=\varnothing,
\]
both closed.  Now fix $c\ge0$.  By local non-triviality, there exists a binary prior $\theta=(s,1-s)$ with
$u:=H(\theta)>0$.  Choose $n\in\mathbb{N}$ and $\lambda\in[0,1]$ such that $c=\lambda nu$
(take $n$ large enough that $c\le nu$, then set $\lambda:=c/(nu)$).

Let $w_n:=\theta^{\otimes n}$ be the $n$-fold product prior on $B_n:=\{0,1\}^n$.  By product additivity,
$H(w_n)=nH(\theta)=nu$.  Define the erasure experiment $E_{n,\lambda}:B_n\to\Delta(\{*,1,\dots,2^n\})$ which
reveals the state fully with probability $\lambda$ and reveals nothing with probability $1-\lambda$.  By the
entropy-reduction formula,
\[
\widehat F_{w_n}(\mu_{w_n,E_{n,\lambda}})=\lambda H(w_n)=\lambda nu=c.
\]
Now consider the block action space $A_c:=\{0,1\}\sqcup B_n$ and the block-diagonal channel
$K_c:=\Id_{\{0,1\}}\sqcup E_{n,\lambda}$.  For $t\in[0,1]$, let $q_t^0$ be the prior $(t,1-t)$ supported on
block $\{0,1\}$, and let $w_n^1$ be $w_n$ supported on block $B_n$.  By the scaled form of
Lemma~\ref{pt:lem:blockbridge} (which extends to boundary priors via Lemma~\ref{pt:lem:supprestrict}),
\[
q_t^0\succeq_{K_c}w_n^1\quad\Longleftrightarrow\quad V(q_t,\Id_{\{0,1\}})=h(t)\ge c=V(w_n,E_{n,\lambda}).
\]
Since $K_c$ is fixed and the map $t\mapsto q_t^0$ is continuous in $\ell_1$, Condition~\textup{(P2)} implies that the
set $\{t\in[0,1]:h(t)\ge c\}$ is closed.  Similarly, $\{t\in[0,1]:h(t)\le c\}$ is closed.  Thus all real superlevel and
sublevel sets of $h$ are closed, so $h$ is both upper and lower semicontinuous on $[0,1]$, hence continuous.

\emph{Continuity on higher simplices by induction.}
For $n\ge3$, write any $q\in\Delta_n$ as $q=(q_1,(1-q_1)\bar q)$ where
$\bar q:=(q_2/(1-q_1),\dots,q_n/(1-q_1))\in\Delta_{n-1}$ for $q_1<1$.  The Faddeev recursion with the
binary partition $\{\{1\},\{2,\dots,n\}\}$ gives
\[
H(q)=h(q_1)+(1-q_1)H(\bar q).
\]
On the region $0<q_1<1$, this is continuous: $h$ is continuous by binary continuity, $H|_{\Delta_{n-1}}$ is
continuous by the induction hypothesis, and $q\mapsto(q_1,\bar q)$ is continuous.  If $q_1=0$, support restriction,
equivalently the positive-mass-block convention in the recursion, gives $H(q)=H(\bar q)$ and $h(0)=0$,
so continuity at that face follows from the induction hypothesis.  At $q_1=1$, the induction hypothesis and compactness
of $\Delta_{n-1}$ imply $H$ is bounded on $\Delta_{n-1}$, so
$(1-q_1)H(\bar q)\to0$ as $q_1\to1$.  Thus $H(q)\to h(1)=0=H(\delta_1)$.  The other boundary faces are handled
by support restriction, relabelling on the support, and pulling back.  This completes the induction.

We have verified the hypotheses of the strong-additivity form of Faddeev's finite entropy characterisation, as stated
for example by Baez--Fritz--Leinster \cite[Theorem~6]{BaezFritzLeinster2011}, following Faddeev's original theorem
\cite{Faddeev1956}: a function on finite probability vectors that is continuous on each closed simplex, invariant under
bijections, expansive under adjoining zero-probability states, zero on point masses, and strongly additive is a
nonnegative multiple of Shannon entropy.  Here invariance under bijections follows from
Corollary~\ref{pt:cor:permutationinvariance}, with boundary cases obtained by restricting to the support and pulling back;
expansibility follows from support restriction; and the grouping recursion above is precisely strong additivity, read
over positive-mass blocks.  Faddeev's theorem therefore gives $H(q)=\alpha\Sh(q)$ for some $\alpha\ge0$.
Local non-triviality (Condition~\textup{(P1)}) forces $\alpha>0$.

\emph{Mutual information.}  Substituting $H=\alpha\Sh$ into the entropy-reduction formula yields
\[
\widehat F_q(\mu_{q,P})
=\alpha\Sh(q)-\sum_{x:m_{q,P}(x)>0}m_{q,P}(x)\alpha\Sh(r_x^{\,q,P})
=\alpha\bigl(\Sh(q)-\E[\Sh(r_X)]\bigr)
=\alpha I_{q,P}(A;X).
\qedhere
\]
\end{proof}

\begin{lemma}[Fixed-environment mutual-information representation]\label{pt:lem:globalsketch}
For every finite channel $P:A\to\Delta(X)$ and all $q,q'\in\Delta(A)$,
\[
q\succeq_P q'
\quad\Longleftrightarrow\quad
I_{q,P}(A;X)\ge I_{q',P}(A;X).
\]
\end{lemma}

\begin{proof}
Let
\[
V(s,R):=\widehat F_s(\mu_{s,R}),
\]
with boundary priors interpreted by Lemma~\ref{pt:lem:supprestrict}.  Lemma~\ref{pt:lem:blockbridge}, divided by the universal
scale from Lemma~\ref{pt:lem:scalecoherence}, gives the cross-prior block-comparison bridge: for finite action sets,
priors $s,t$, and channels $R,S$,
\[
s^0\succeq_{R\sqcup S}t^1
\quad\Longleftrightarrow\quad
V(s,R)\ge V(t,S),
\]
including boundary priors by support restriction.  Lemma~\ref{pt:lem:faddeevsketch} gives, for every finite
prior/channel pair,
\[
V(s,R)=\alpha I_{s,R}
\]
with $\alpha>0$.

Fix $P:A\to\Delta(X)$ and $q,q'\in\Delta(A)$.  By the duplication clause of Condition~\textup{(P3)},
\[
q\succeq_P q'
\quad\Longleftrightarrow\quad
q^0\succeq_{P\sqcup P}(q')^1.
\]
Applying this bridge with $(s,R)=(q,P)$ and $(t,S)=(q',P)$ gives
\[
q^0\succeq_{P\sqcup P}(q')^1
\quad\Longleftrightarrow\quad
V(q,P)\ge V(q',P).
\]
Substituting $V=\alpha I$ and using $\alpha>0$ yields
\[
q\succeq_P q'
\quad\Longleftrightarrow\quad
I_{q,P}(A;X)\ge I_{q',P}(A;X).
\]

If $q$ or $q'$ is on the boundary, the bridge is read after restricting each pair to its positive-probability support
and then pulling the comparison back by Lemma~\ref{pt:lem:supprestrict}.  Deleting
zero-probability actions does not change mutual information, and the Condition~\textup{(P3)} equivalence remains the ambient
fixed-channel comparison for $P$.
\end{proof}

\begin{proof}[Proof of Proposition~\ref{pt:prop:main}]
That \textup{(ii)} implies \textup{(i)} is the necessity verification at the start of this appendix.  That
\textup{(i)} implies \textup{(ii)} is Lemma~\ref{pt:lem:globalsketch}.  For the moreover clause, apply \textup{(ii)} to
the fixed block environment $\bigsqcup_{k\in K}P_k$: a lottery supported on block $i$ has constant block label, so
\[
I_{q_i^i,\bigsqcup_{k\in K}P_k}=I_{q_i,P_i}(A_i;X_i),
\]
and similarly for block $j$.
\end{proof}


\begin{thebibliography}{9}

\bibitem{BaezFritzLeinster2011}
John C. Baez, Tobias Fritz, and Tom Leinster.
\newblock A characterization of entropy in terms of information loss.
\newblock {\em Entropy}, 13(11):1945--1957, 2011.

\bibitem{Faddeev1956}
D. K. Faddeev.
\newblock On the concept of entropy of a finite probabilistic scheme.
\newblock {\em Uspekhi Matematicheskikh Nauk}, 11(1):227--231, 1956.
\newblock In Russian.

\bibitem{HersteinMilnor1953}
I. N. Herstein and John Milnor.
\newblock An axiomatic approach to measurable utility.
\newblock {\em Econometrica}, 21(2):291--297, 1953.

\end{thebibliography}

\end{document}
